\documentclass[10pt,reqno]{amsart}

\usepackage[T1]{fontenc}
\usepackage{amsmath,amssymb}
\usepackage{microtype}
\usepackage[margin=1in]{geometry}
\usepackage{calc}
\usepackage{booktabs,longtable,array}
\usepackage{enumitem}
\usepackage[hidelinks]{hyperref}
\usepackage[nameinlink,capitalise,noabbrev]{cleveref}
\pdfstringdefDisableCommands{%
  \def\({}\def\){}%
  \def\frac#1#2{#1/#2}%
  \def\ge{>=}\def\le{<=}%
  \def\Delta{Delta}\def\sigma{sigma}\def\nu{nu}%
}
\hypersetup{
  pdftitle={A Bellman-function bound for the Beurling--Ahlfors transform},
  pdfauthor={Hyra},
  pdfsubject={An explicit uniform Lp bound for the Beurling--Ahlfors transform},
  pdfkeywords={Beurling--Ahlfors transform, Bellman function, conformal martingale}
}

\allowdisplaybreaks
\emergencystretch=3em
\providecommand{\tightlist}{\setlength{\itemsep}{0pt}\setlength{\parskip}{0pt}}
\newcommand{\C}{\mathbb C}
\newcommand{\R}{\mathbb R}
\newcommand{\E}{\mathbb E}
\newcommand{\pv}{\operatorname{p.v.}}
\newcommand{\artanh}{\operatorname{artanh}}

\newtheorem{theorem}{Theorem}[section]
\newtheorem{lemma}[theorem]{Lemma}
\newtheorem{proposition}[theorem]{Proposition}
\theoremstyle{definition}
\newtheorem{definition}[theorem]{Definition}
\theoremstyle{remark}
\newtheorem{remark}[theorem]{Remark}
\newtheorem*{maintheorem}{Main Theorem}
\newtheorem*{bellmanprinciple}{Bellman principle}
\newtheorem*{boundsummary}{Endpoint theorem}

\title[A Bellman-function bound for the Beurling--Ahlfors transform]
{A Bellman-function bound for the Beurling--Ahlfors transform}
\author{Hyra}
\date{August 18, 2026}
\subjclass[2020]{Primary 42B20; Secondary 60G44}
\keywords{Beurling--Ahlfors transform, Bellman function, conformal martingale,
heat extension, singular integral, Riesz--Thorin interpolation}

\begin{document}

\begin{abstract}
We prove the uniform estimate
\[
 \|B\|_{L^p(\mathbb C)\to L^p(\mathbb C)}
 \le 1.523958\,(p^*-1),\qquad 1<p<\infty,
\]
for the Beurling--Ahlfors transform.  The argument combines a heat-extension
representation with an explicit Bellman-function criterion adapted to the
conformal martingale structure of the transform.  The principal new ingredients
are a two-sided kinematic constraint that reduces the drift calculation to one
negative-semidefinite \(2\times2\) matrix, and two explicit generalized
polynomial Bellman profiles at \(p=16/3\) and \(p=15/2\).  Their drift and
majorization properties are certified by exact polynomial identities and rational
inequalities.  Interpolation with the \(L^2\) isometry and a self-contained
Burkholder-type bound yields the stated constant.
\end{abstract}

\maketitle

\section*{Introduction}
The Beurling--Ahlfors transform is the singular integral
\[
 Bf(z)=-\frac1\pi\,\pv\!\int_{\mathbb C}
       \frac{f(w)}{(z-w)^2}\,dm(w).
\]
For \(1<p<\infty\), write
\(p^*=\max\{p,p/(p-1)\}\).  Iwaniec's conjecture asserts that
\(\|B\|_p=p^*-1\); see \cite{Iwaniec82}.  Bañuelos and Janakiraman proved the
uniform upper bound \(1.575(p^*-1)\) by exploiting conformal martingales
\cite{BanuelosJanakiraman08}.  The present paper lowers that coefficient.

\begin{maintheorem}\label{thm:main}
For every \(1<p<\infty\),
\[
 \|B\|_{L^p(\mathbb C)\to L^p(\mathbb C)}
 \le 1.523958\,(p^*-1).
\]
Consequently, if
\[
 C_{\mathrm{BA}}:=\sup_{1<p<\infty}
 \frac{\|B\|_{L^p\to L^p}}{p^*-1},
\]
then \(C_{\mathrm{BA}}\le1.523958\).
\end{maintheorem}

The proof follows the heat-extension and space-time Brownian-motion framework
of \cite{BanuelosMendez03}, but no quantitative martingale inequality is imported.
Instead, the necessary martingale estimate is proved from an explicit radial
Bellman function.  The structural step is a two-sided inequality
\(|T-P|\le W\), which permits the angular variable to be eliminated without a
case split.  This leads to a pointwise drift criterion expressed by a single
\(2\times2\) matrix.  Generalized polynomial profiles with real exponents then
produce the endpoint estimates
\[
 \|B\|_{16/3}\le22830^{3/16},\qquad
 \|B\|_{15/2}\le22865000^{2/15}.
\]
Together with \(\|B\|_2=1\), a self-contained estimate
\(\|B\|_p\le\sqrt{2p(p-1)}\) for \(p\ge2\), duality, and
Riesz--Thorin interpolation, these endpoints imply the main theorem.

All profile identities and all numerical comparisons used in the proof are exact.
The only external inputs are standard results from Fourier analysis, functional
analysis, and stochastic calculus, together with the classical multiplier
identification of \(B\).  The use of compactly supported mollifiers below makes
the regularization step entirely local and avoids any hidden tail estimate.

\setcounter{section}{-1}
\section{Notation and conventions}\label{sec:notation}

\begin{itemize}
\tightlist
\item
  \(m\) = Lebesgue measure on \(\mathbb C\simeq\mathbb R^2\); we write \(dx\) for it and identify \(z=x_1+ix_2\leftrightarrow x=(x_1,x_2)\). All functions are complex valued unless stated otherwise. \(\|f\|_q=(\int_{\mathbb R^2}|f|^q\,dx)^{1/q}\).
\item
  Fourier transform: \(\hat f(\xi)=\int_{\mathbb R^2}f(x)e^{-2\pi i x\cdot\xi}\,dx\), so that Plancherel reads \(\int f\bar g\,dx=\int\hat f\overline{\hat g}\,d\xi\) and \(\widehat{\partial_j f}(\xi)=2\pi i\xi_j\hat f(\xi)\). We also identify \(\xi=\xi_1+i\xi_2\in\mathbb C\) when convenient; \(\bar\xi=\xi_1-i\xi_2\), \(|\xi|^2=\xi_1^2+\xi_2^2\).
\item
  Heat semigroup: \(G_t(x)=(4\pi t)^{-1}e^{-|x|^2/(4t)}\) for \(t>0\) (a probability density on \(\mathbb R^2\)), and for \(f\in C_c^\infty(\mathbb R^2)\) \[U(x,t):=(e^{t\Delta}f)(x)=(f*G_t)(x),\qquad U(x,0)=f(x).\] Then \(\partial_tU=\Delta U\), \(\hat U(\xi,t)=e^{-4\pi^2|\xi|^2t}\hat f(\xi)\), \(U\in C^\infty(\mathbb R^2\times[0,\infty))\) (differentiation under the integral sign is legitimate because \(f\in C_c^\infty\) and \(\partial^\alpha U=(\partial^\alpha f)*G_t\)), and by Young's inequality \(\|\partial^\alpha U(\cdot,t)\|_q\le\|\partial^\alpha f\|_q\) for every multi-index \(\alpha\), every \(q\in[1,\infty]\) and every \(t\ge0\).
\item
  For \(v,w\in\mathbb C^2\) we write \(v\cdot w=\sum_j v_jw_j\) (\textbf{bilinear}, no conjugation) and \(|v|^2=\sum_j|v_j|^2\); thus \(|v\cdot w|\le|v|\,|w|\). For a real vector \(h\in\mathbb R^d\), \(|h|\) is the Euclidean norm and \(h\cdot k\) the real inner product.
\item
  \(\mathbb C\) is identified with \(\mathbb R^2\) as a real Euclidean space whenever a complex quantity is fed into a real function; the modulus of a complex number equals the Euclidean norm of the corresponding real vector, and for \(\zeta,\eta\in\mathbb C\) the real inner product is \(\zeta\cdot\eta=\operatorname{Re}(\zeta\bar\eta)\).
\item
  \(\operatorname{artanh}x=\frac12\ln\frac{1+x}{1-x}\) for \(|x|<1\).
\end{itemize}

\paragraph{The multiplier form of \(B\).}
We use the classical identity
\[\widehat{Bf}(\xi)=\frac{\bar\xi}{\xi}\,\widehat f(\xi)
 =\frac{\bar\xi^{\,2}}{|\xi|^2}\widehat f(\xi),
 \qquad \xi\ne0, \tag{0.1}\]
initially for \(f\in C_c^\infty(\mathbb R^2)\).  In particular, Plancherel and
\(|\bar\xi/\xi|=1\) give
\[\|B\|_{L^2\to L^2}=1. \tag{0.2}\]
All other \(L^p\)-extensions used below are obtained from the a priori estimates
proved in this paper.  More precisely, an estimate on \(C_c^\infty\) gives a
unique bounded extension by density.  If bounds are available at two exponents,
the corresponding extensions agree on their intersection because
\(C_c^\infty\) is dense there for the sum of the two norms.  Thus the compatible
form of the Riesz--Thorin theorem applies without assuming qualitative
\(L^q\)-boundedness in advance.

\section{Reductions and the heat representation}\label{sec:heat}

\subsection{Reduction to \(p\ge2\)}\label{subsec:duality}

\begin{lemma}\label{lem:1-1}
For \(f,g\in C_c^\infty(\mathbb R^2)\) one has \(\int (Bf)g\,dx=\int f\,(Bg)\,dx\).
\end{lemma}

\begin{proof}
For \(F,G\in L^2\), Plancherel in the bilinear form reads \(\int FG\,dx=\int\hat F(\xi)\hat G(-\xi)\,d\xi\) (apply \(\int F\bar H=\int\hat F\overline{\hat H}\) with \(H=\bar G\), whose transform is \(\hat H(\xi)=\overline{\hat G(-\xi)}\)). Hence, with \(m(\xi)=\bar\xi/\xi\), \[\int (Bf)g\,dx=\int m(\xi)\hat f(\xi)\hat g(-\xi)\,d\xi,\qquad
\int f\,(Bg)\,dx=\int \hat f(\xi)\,m(-\xi)\hat g(-\xi)\,d\xi .\] Since \(m(-\xi)=\overline{(-\xi)}/(-\xi)=\bar\xi/\xi=m(\xi)\), the two agree.
\end{proof}

\begin{lemma}[duality]\label{lem:1-2}
Let \(1<q<\infty\), \(q'=q/(q-1)\), let \(h\in L^1_{\mathrm{loc}}(\mathbb R^2)\) and \(M<\infty\) be such that \(\big|\int h\bar g\,dx\big|\le M\|g\|_{q}\) for every \(g\in C_c^\infty\). Then \(h\in L^{q'}\) and \(\|h\|_{q'}\le M\).
\end{lemma}

\begin{proof}
The map \(g\mapsto\int h\bar g\,dx\) is a conjugate-linear functional on the dense subspace \(C_c^\infty\) of \(L^{q}\), of norm \(\le M\); extend it continuously to \(L^{q}\). Since \(1<q<\infty\), the dual of \(L^{q}\) is \(L^{q'}\), so there is \(H\in L^{q'}\) with \(\|H\|_{q'}\le M\) and \(\int h\bar g=\int H\bar g\) for every \(g\in C_c^\infty\). Then \(h-H\in L^1_{\mathrm{loc}}\) annihilates \(C_c^\infty\), hence \(h=H\) a.e., and \(\|h\|_{q'}=\|H\|_{q'}\le M\).
\end{proof}

\begin{lemma}[transposition]\label{lem:1-3}
Suppose \(\|Bg\|_{q}\le M_q\|g\|_q\) for all \(g\in C_c^\infty\) and some \(1<q<\infty\). Then \(\|Bf\|_{q'}\le M_q\|f\|_{q'}\) for all \(f\in C_c^\infty\), where \(1/q+1/q'=1\). Consequently \(\|B\|_{L^p\to L^p}=\|B\|_{L^{p'}\to L^{p'}}\) for all \(1<p<\infty\).
\end{lemma}

\begin{proof}
Let \(f\in C_c^\infty\); then \(Bf\in L^2\subset L^1_{\mathrm{loc}}\). For every \(g\in C_c^\infty\) we also have \(\bar g\in C_c^\infty\), so Lemma 1.1 (applied to the pair \(f,\bar g\)) and Hölder give \[\Big|\int (Bf)\bar g\,dx\Big|=\Big|\int f\,(B\bar g)\,dx\Big|\le\|f\|_{q'}\|B\bar g\|_q\le M_q\|f\|_{q'}\|g\|_q .\] Lemma 1.2 (with \(h=Bf\), \(M=M_q\|f\|_{q'}\)) gives \(\|Bf\|_{q'}\le M_q\|f\|_{q'}\). By density \(\|B\|_{q'}\le\|B\|_q\), and by symmetry (replace \(q\) by \(q'\)) the reverse inequality holds as well.
\end{proof}

Since \(p^*-1=q-1\) both for \(p=q\ge 2\) and for \(p=q'\le2\), Lemma 1.3 shows:

\begin{quote}
\textbf{It suffices to prove \(\|Bf\|_p\le\Lambda(p)\|f\|_p\) for all \(f\in C_c^\infty\) and all \(p\ge2\)}, where \(\Lambda(p)\) is the desired bound; the case \(1<p<2\) then follows with \(\Lambda(p')\).
\end{quote}

For \(p\ge 2\) we write \(p'=p/(p-1)\in(1,2]\) and \(q:=p-1\ge1\).

\subsection{The matrix \(A\) and conformality}\label{subsec:matrix}

Let \[w:=(1,-i)\in\mathbb C^2,\qquad A:=w\,w^{T}=\begin{pmatrix}1&-i\\-i&-1\end{pmatrix}
\in\mathbb C^{2\times2}.\] For a real vector \(\xi=(\xi_1,\xi_2)\), \[\langle A\xi,\xi\rangle:=\sum_{j,k}A_{jk}\xi_j\xi_k=(w\cdot\xi)^2=(\xi_1-i\xi_2)^2=\bar\xi^{\,2},\] so by (0.1) the symbol of \(B\) is exactly \(\langle A\xi,\xi\rangle/|\xi|^2\).

\begin{lemma}[structure of \(A\)]\label{lem:1-4}
Let \(v=(v_1,v_2)\in\mathbb C^2\) and put \(c:=w\cdot v=v_1-iv_2\in\mathbb C\). Then \[Av=(c,\,-ic),\qquad |Av|^2=2|c|^2\le 4|v|^2 . \tag{1.1}\] Moreover, identifying \(\mathbb C\simeq\mathbb R^2\), the two components \(k_1:=c\) and \(k_2:=-ic\) of \(Av\) satisfy, \textbf{for every unit vector \(\hat e\in\mathbb R^2\)}, \[\sum_{j=1}^2(k_j\cdot\hat e)^2=|c|^2=\tfrac12\sum_{j=1}^2|k_j|^2,\qquad
\sum_{j=1}^2\big|k_j-(k_j\cdot\hat e)\hat e\big|^2=|c|^2 . \tag{1.2}\]
\end{lemma}

\begin{proof}
\(Av=w(w\cdot v)=cw=(c,-ic)\), and \(|Av|^2=|c|^2+|-ic|^2=2|c|^2\); also \(|c|=|v_1-iv_2|\le|v_1|+|v_2|\le\sqrt2|v|\), whence \(|Av|^2\le4|v|^2\). For (1.2): write \(\hat e\in\mathbb R^2\) as a unit complex number \(e\), so that \(k\cdot\hat e=\operatorname{Re}(k\bar e)\) for \(k\in\mathbb C\). Then \[(k_1\cdot\hat e)^2+(k_2\cdot\hat e)^2
=\big(\operatorname{Re}(c\bar e)\big)^2+\big(\operatorname{Re}(-ic\bar e)\big)^2
=\big(\operatorname{Re}(c\bar e)\big)^2+\big(\operatorname{Im}(c\bar e)\big)^2=|c\bar e|^2=|c|^2 ,\] using \(\operatorname{Re}(-i\zeta)=\operatorname{Im}(\zeta)\) and \(|e|=1\). Since \(|k_j|^2=|c|^2\) for \(j=1,2\), we get \(\sum_j|k_j|^2=2|c|^2\) and \(\sum_j|k_j-(k_j\cdot\hat e)\hat e|^2=\sum_j\big(|k_j|^2-(k_j\cdot\hat e)^2\big)=2|c|^2-|c|^2=|c|^2\).
\end{proof}

The identity (1.2) expresses the fact that \(\{k_1,k_2\}\) is an \emph{orthogonal pair of equal length}: \(k_2=-ik_1\) is \(k_1\) rotated by \(-\pi/2\). Any such pair distributes its total energy \(2|c|^2\) \textbf{isotropically}: exactly half of it lies along any prescribed direction and half orthogonally to it. The inequality \(|Av|\le2|v|\) in (1.1) will be used only qualitatively (to know that certain \(L^2\)-quantities are finite); the \emph{quantitative} input is (1.2).

\subsection{Heat representation of the bilinear form}\label{subsec:heat-rep}

\begin{lemma}\label{lem:1-5}
Let \(f,g\in C_c^\infty(\mathbb R^2)\), \(U=e^{t\Delta}f\), \(V=e^{t\Delta}g\). Then \[\int_{\mathbb R^2}(Bf)\,\bar g\,dx\;=\;2\int_0^\infty\!\!\int_{\mathbb R^2}
\big(A\nabla U(x,t)\big)\cdot\overline{\nabla V(x,t)}\;dx\,dt , \tag{1.3}\] where \(\nabla U=(\partial_1U,\partial_2U)\in\mathbb C^2\) and the dot is the bilinear pairing of Section~0 (so the right side is \(\sum_{j,k}A_{jk}\int\!\!\int\partial_kU\,\overline{\partial_jV}\)). Moreover \(\int_0^\infty\!\int|\nabla U||\nabla V|\,dx\,dt<\infty\) and \[\int_0^\infty\!\!\int_{\mathbb R^2}|\nabla U|^2\,dx\,dt=\tfrac12\|f\|_2^2 . \tag{1.4}\]
\end{lemma}

\begin{proof}
Fix \(t>0\). By Plancherel, for each \(j,k\), \[\int \partial_kU\,\overline{\partial_jV}\,dx=\int (2\pi i\xi_k\hat U)\overline{(2\pi i\xi_j\hat V)}\,d\xi
=4\pi^2\int \xi_j\xi_k\,e^{-8\pi^2|\xi|^2t}\hat f\overline{\hat g}\,d\xi .\] Summing against \(A_{jk}\), \[\int (A\nabla U)\cdot\overline{\nabla V}\,dx=4\pi^2\int\langle A\xi,\xi\rangle e^{-8\pi^2|\xi|^2t}\hat f\overline{\hat g}\,d\xi .\] The double integral converges absolutely: \(|\langle A\xi,\xi\rangle|=|\bar\xi^2|=|\xi|^2\) and \(\int_0^\infty\!\int 4\pi^2|\xi|^2e^{-8\pi^2|\xi|^2t}|\hat f||\hat g|\,d\xi\,dt
=\frac12\int|\hat f||\hat g|\,d\xi<\infty\) (Tonelli; \(\hat f,\hat g\) are Schwartz). Hence by Fubini and \(\int_0^\infty e^{-8\pi^2|\xi|^2t}dt=(8\pi^2|\xi|^2)^{-1}\), \[2\int_0^\infty\!\!\int (A\nabla U)\cdot\overline{\nabla V}\,dx\,dt
=2\cdot\frac{4\pi^2}{8\pi^2}\int\frac{\langle A\xi,\xi\rangle}{|\xi|^2}\hat f\overline{\hat g}\,d\xi
=\int\frac{\bar\xi^{\,2}}{|\xi|^2}\hat f\overline{\hat g}\,d\xi=\int(Bf)\bar g\,dx,\] the last step by (0.1) and Plancherel. For the final assertions, Plancherel gives \[\int_0^\infty\!\!\int|\nabla U|^2dx\,dt=\int_0^\infty\!\!\int4\pi^2|\xi|^2e^{-8\pi^2|\xi|^2t}|\hat f|^2d\xi\, dt=\tfrac12\|f\|_2^2<\infty,\] and similarly for \(V\); then Cauchy--Schwarz yields \(\int_0^\infty\!\int|\nabla U||\nabla V|\le\frac12\|f\|_2\|g\|_2<\infty\).
\end{proof}

Thus everything is reduced to estimates of the form \[\Big|2\int_0^\infty\!\!\int (A\nabla U)\cdot\overline{\nabla V}\,dx\,dt\Big|
\ \le\ \Lambda\,\|f\|_p\|g\|_{p'} ,\qquad p\ge2, \tag{1.5}\] which are proved in Section~3--Section~4.


\section{Admissible Bellman functions}\label{sec:bellman}

Throughout this section \(p\ge2\) is a fixed real number.

\subsection{The two-sided kinematic constraint}\label{subsec:kinematic}

\begin{lemma}\label{lem:2-1}
Let \(a,b\in\mathbb C\), put \(c:=a-ib\) and \(T:=|c|\), and let \(\hat e\in\mathbb R^2\) be a unit vector. Set \(h_1:=a\), \(h_2:=b\) (viewed in \(\mathbb R^2\)) and \[P:=\Big(\sum_{j=1}^2(h_j\cdot\hat e)^2\Big)^{1/2},\qquad
W:=\Big(\sum_{j=1}^2\big|h_j-(h_j\cdot\hat e)\hat e\big|^2\Big)^{1/2} .\] Then \[|\,T-P\,|\ \le\ W ,\qquad\text{equivalently}\qquad T\le P+W\ \text{ and }\ P\le T+W. \tag{2.1}\]
\end{lemma}

\begin{proof}
Let \(\varepsilon\in\mathbb C\), \(|\varepsilon|=1\), correspond to \(\hat e\), and set \(a':=a\bar\varepsilon=a_1+ia_2\), \(b':=b\bar\varepsilon=b_1+ib_2\) with \(a_i,b_i\in\mathbb R\). Since \(h\cdot\hat e=\operatorname{Re}(h\bar\varepsilon)\) and \(|h-(h\cdot\hat e)\hat e|=|\operatorname{Im}(h\bar\varepsilon)|\) for \(h\in\mathbb C\), \[P^2=a_1^2+b_1^2,\qquad W^2=a_2^2+b_2^2 .\] Moreover \(c\bar\varepsilon=a'-ib'=(a_1+b_2)+i(a_2-b_1)\), so, as vectors in \(\mathbb R^2\), \[T=|c|=|c\bar\varepsilon|=\big|(a_1+b_2,\;a_2-b_1)\big|=\big|\,\mathsf u+\mathsf v\,\big|,\qquad
\mathsf u:=(a_1,-b_1),\quad \mathsf v:=(b_2,a_2),\] and \(|\mathsf u|=P\), \(|\mathsf v|=W\). The triangle inequality gives \(T=|\mathsf u+\mathsf v|\le|\mathsf u|+|\mathsf v|=P+W\) and \(P=|\mathsf u|=|(\mathsf u+\mathsf v)-\mathsf v|\le T+W\).
\end{proof}

\begin{quote}
Only the first half of (2.1) is classical. The second half, \(P\le T+W\), is what allows us below to replace \(W^2\) by \((T-P)^2\) \textbf{without a case distinction}, and hence to reduce the whole drift condition to the negative semidefiniteness of a single \(2\times2\) matrix.
\end{quote}

\subsection{The drift criterion}\label{subsec:drift}

Let \(\varphi:[0,\infty)^2\to\mathbb R\) be continuous, and set \[u(x,y):=\varphi(|x|,|y|),\qquad x,y\in\mathbb R^2 .\] Suppose \(\varphi\) is \(C^2\) on the open quadrant \(\{r>0,\ s>0\}\). There, with \(\hat x=x/|x|\), \(\hat y=y/|y|\), \(r=|x|\), \(s=|y|\), the classical Hessian of \(u\) is \[D^2_{xx}u=\varphi_{rr}\,\hat x\otimes\hat x+\frac{\varphi_r}{r}(I-\hat x\otimes\hat x),\quad
D^2_{xy}u=\varphi_{rs}\,\hat x\otimes\hat y,\quad
D^2_{yy}u=\varphi_{ss}\,\hat y\otimes\hat y+\frac{\varphi_s}{s}(I-\hat y\otimes\hat y)\] (these follow from \(\partial_{x_i}|x|=\hat x_i\) and \(\partial_{x_i}\hat x_j=(\delta_{ij}-\hat x_i\hat x_j)/r\), so that \(\partial_{x_i}\partial_{x_j}u=\varphi_{rr}\hat x_i\hat x_j+\varphi_r(\delta_{ij}-\hat x_i\hat x_j)/r\), and similarly in \(y\); the mixed block is \(\partial_{x_i}\partial_{y_j}u=\varphi_{rs}\hat x_i\hat y_j\)). Hence for \(h,k\in\mathbb R^2\), writing \(h_{\|}:=h\cdot\hat x\), \(h_\perp:=h-h_{\|}\hat x\), \(k_{\|}:=k\cdot\hat y\), \(k_\perp:=k-k_{\|}\hat y\), \[\big\langle D^2u(x,y)(h,k),(h,k)\big\rangle
=\varphi_{rr}h_{\|}^2+2\varphi_{rs}h_{\|}k_{\|}+\varphi_{ss}k_{\|}^2
+\frac{\varphi_r}{r}|h_\perp|^2+\frac{\varphi_s}{s}|k_\perp|^2. \tag{2.2}\]

\begin{definition}[the drift matrix]\label{def:drift-matrix}
For \(r,s>0\) put \[\mathcal A:=\varphi_{rr}+\frac{\varphi_r}{r},\qquad
\mathcal B:=|\varphi_{rs}|-\frac{\varphi_r}{r},\qquad
\mathcal C:=\varphi_{ss}+\frac{\varphi_s}{s}+\frac{\varphi_r}{r},\] all evaluated at \((r,s)\), and call \(\mathcal Q:=\begin{pmatrix}\mathcal A&\mathcal B\\ \mathcal B&\mathcal C\end{pmatrix}\) the \emph{drift matrix} of \(\varphi\) at \((r,s)\).
\end{definition}

\begin{proposition}[drift inequality]\label{prop:2-3}
Let \(x,y\in\mathbb R^2\) with \(x\ne0\ne y\), and let \(a,b\in\mathbb C\) be arbitrary. Put \(c:=a-ib\) and \[v_1:=(a,\,c),\qquad v_2:=(b,\,-ic)\ \in\ \mathbb R^2\times\mathbb R^2=\mathbb R^4 .\] Assume that at \((r,s)=(|x|,|y|)\) one has \[\varphi_r\le0 \tag{D0}\] and that the drift matrix is negative semidefinite, i.e. \[\mathcal A\le0,\qquad \mathcal C\le0,\qquad \mathcal B^2\le\mathcal A\,\mathcal C. \tag{D1}\] Then \[\sum_{j=1}^{2}\big\langle D^2u(x,y)\,v_j,\,v_j\big\rangle\ \le\ 0 . \tag{2.3}\]
\end{proposition}

\begin{proof}
Write \(h_1=a,\ h_2=b,\ k_1=c,\ k_2=-ic\); the \(j\)-th vector is \(v_j=(h_j,k_j)\). Apply (2.2) to each \(j\) and sum. Let \(P,W\) be as in Lemma 2.1 with \(\hat e=\hat x\), so \(\sum_jh_{j\|}^2=P^2\) and \(\sum_j|h_{j\perp}|^2=W^2\); and let \(T=|c|\). Lemma 1.4, applied with \(\hat e=\hat y\), gives \[\sum_j k_{j\|}^2=T^2,\qquad \sum_j|k_{j\perp}|^2=T^2 .\] By Cauchy--Schwarz, \(\big|\sum_jh_{j\|}k_{j\|}\big|\le(\sum_jh_{j\|}^2)^{1/2}(\sum_jk_{j\|}^2)^{1/2}=PT\). Therefore \[\sum_{j}\big\langle D^2u\,v_j,v_j\big\rangle
\ \le\ \varphi_{rr}P^2+2|\varphi_{rs}|\,PT+\Big(\varphi_{ss}+\frac{\varphi_s}{s}\Big)T^2
+\frac{\varphi_r}{r}\,W^2 .\] By (D0) the coefficient \(\varphi_r/r\) of \(W^2\) is \(\le0\), and by Lemma 2.1 \(W^2\ge(T-P)^2\); hence \(\frac{\varphi_r}{r}W^2\le\frac{\varphi_r}{r}(T-P)^2\) and \[\sum_{j}\big\langle D^2u\,v_j,v_j\big\rangle
\ \le\ \Big(\varphi_{rr}+\frac{\varphi_r}{r}\Big)P^2
+2\Big(|\varphi_{rs}|-\frac{\varphi_r}{r}\Big)PT
+\Big(\varphi_{ss}+\frac{\varphi_s}{s}+\frac{\varphi_r}{r}\Big)T^2
=\mathcal AP^2+2\mathcal BPT+\mathcal CT^2 .\] A real binary quadratic form \(\mathcal AP^2+2\mathcal BPT+\mathcal CT^2\) with \(\mathcal A\le0\), \(\mathcal C\le0\) and \(\mathcal B^2\le\mathcal A\mathcal C\) is negative semidefinite on all of \(\mathbb R^2\): indeed if \(\mathcal A<0\) then \(\mathcal AP^2+2\mathcal BPT+\mathcal CT^2
=\mathcal A\big(P+\frac{\mathcal B}{\mathcal A}T\big)^2
+\frac{\mathcal A\mathcal C-\mathcal B^2}{\mathcal A}T^2\le0\), while if \(\mathcal A=0\) then \(\mathcal B^2\le0\) forces \(\mathcal B=0\) and the form is \(\mathcal CT^2\le0\). This proves (2.3).
\end{proof}

\subsection{Admissible functions}\label{subsec:admissible}

\begin{definition}\label{def:admissible}
Let \(p\ge2\) and \(\Lambda>0\). A function \(\varphi:[0,\infty)^2\to\mathbb R\) is called \textbf{\((p,\Lambda)\)-admissible} if:

\begin{itemize}
\tightlist
\item
  \textbf{(A1)} \(\varphi\) is continuous on \([0,\infty)^2\), positively homogeneous of degree \(p\) (i.e.~\(\varphi(\lambda r,\lambda s)=\lambda^p\varphi(r,s)\) for \(\lambda>0\)), and \(C^\infty\) on \(\{r>0,s>0\}\);
\item
  \textbf{(A2)} \emph{(growth)} there are \(M<\infty\) and \(\theta\in[0,1)\) with, for all \(r,s>0\) and \(\sigma:=r+s\), \[|\varphi|\le M\sigma^{p},\qquad
  |\varphi_r|+|\varphi_s|\le M\sigma^{p-1}\Big(\frac\sigma s\Big)^{\theta},\] \[|\varphi_{rr}|+|\varphi_{rs}|+|\varphi_{ss}|+\Big|\frac{\varphi_r}{r}\Big|
  +\Big|\frac{\varphi_s}{s}\Big|\ \le\ M\sigma^{p-2}\Big(\frac\sigma s\Big)^{1+\theta};\]
\item
  \textbf{(A3)} \emph{(drift)} (D0) and (D1) of Proposition 2.3 hold at every \((r,s)\) with \(r,s>0\);
\item
  \textbf{(A4)} \emph{(majorization)} \(\varphi(r,s)\ \ge\ s^p-\Lambda^p r^p\) for all \(r,s\ge0\);
\item
  \textbf{(A5)} \emph{(boundary sign)} \(\varphi(r,0)\le0\) for all \(r\ge0\).
\end{itemize}
\end{definition}

Note \(\sigma/s=1+r/s\ge1\), so (A2) with \(\theta=0\) is the strongest form and reads \(|\varphi_r|+|\varphi_s|\le M\sigma^{p-1}\) and (second order) \(\le M\sigma^{p-2}(1+r/s)\); larger \(\theta<1\) makes (A2) weaker. The single purpose of \(\theta\) is to allow profiles whose \(s\)-derivatives blow up like \(s^{-\theta}\) at \(s=0\); everything in Section~2.4 and Section~3 goes through unchanged for any \(\theta<1\).

We will prove in Section~3:

\begin{bellmanprinciple}
If \(\varphi\) is \((p,\Lambda)\)-admissible for some \(p\ge2\), then
\[\|B\|_{L^p\to L^p}\le\Lambda.\]
\end{bellmanprinciple}

Three admissible functions are exhibited in Section~4.

\subsection{Regularity and mollification}\label{subsec:mollification}

Let \(\varphi\) be \((p,\Lambda)\)-admissible with growth parameters \((M,\theta)\), and let \(u(x,y)=\varphi(|x|,|y|)\) on \(\mathbb R^4\). We write \(z=(x,y)\) and \(s=|y|\).

\begin{lemma}\label{lem:2-5}
\(u\) is continuous on \(\mathbb R^4\), and for some \(M'=M'(M,p)\), \[|u(z)|\le M'|z|^p,\qquad
|\nabla u(z)|\le M'\,|z|^{p-1}\Big(\frac{|z|}{|y|}\Big)^{\theta}\quad(y\ne0).
\tag{2.4}\] The pointwise gradient and Hessian of \(u\) exist off the null set \(\mathcal N:=\{x=0\}\cup\{y=0\}\), belong to \(L^1_{\mathrm{loc}}(\mathbb R^4)\), and coincide with the distributional gradient and Hessian of \(u\).
\end{lemma}

\begin{proof}
Continuity and the first bound in (2.4) are immediate from (A1)--(A2) (note \(\sigma=|x|+|y|\le\sqrt2\,|z|\) and \(|z|\le\sigma\)). Off \(\mathcal N\), \(u\) is smooth and \(|\nabla u|\le|\varphi_r|+|\varphi_s|\le M\sigma^{p-1}(\sigma/s)^\theta\), which is the second bound in (2.4).

By (A2), \(|D^2u(z)|\le C_pM\sigma^{p-2}(\sigma/s)^{1+\theta}\) off \(\mathcal N\) (each entry of \(D^2u\) is a combination of \(\varphi_{rr},\varphi_{rs},\varphi_{ss},\varphi_r/r,\varphi_s/s\) with coefficients bounded by \(1\), by the display preceding (2.2)). Hence on the ball \(\{|z|<R\}\), \[|\nabla u(z)|\le C(R)\,|y|^{-\theta},\qquad |D^2u(z)|\le C(R)\,|y|^{-1-\theta}.\] Both are locally integrable on \(\mathbb R^4\): in polar coordinates in the \(y\)-variable, \(\int_{|y|<R}|y|^{-\alpha}dy=2\pi R^{2-\alpha}/(2-\alpha)<\infty\) whenever \(\alpha<2\), and here \(\alpha=\theta<1\) resp. \(\alpha=1+\theta<2\).

Let \(\zeta\in C_c^\infty(\mathbb R^4)\), supported in \(\{|z|<R\}\), and let \(\mathsf a,\mathsf b\) be two of the four coordinates. Let \(\Omega_\delta:=\{(x,y):|x|>\delta,\ |y|>\delta\}\). On \(\Omega_\delta\) the function \(u\) is \(C^\infty\), so two integrations by parts give \[\int_{\Omega_\delta}u\,\partial_{\mathsf a}\partial_{\mathsf b}\zeta
=\int_{\Omega_\delta}(\partial_{\mathsf a}\partial_{\mathsf b}u)\,\zeta+R_\delta,\] where \(R_\delta\) is a sum of boundary integrals over \(\partial\Omega_\delta\subset(\{|x|=\delta\}\times\mathbb R^2)\cup(\mathbb R^2\times\{|y|=\delta\})\) of terms \(u\,(\partial\zeta)\,n\) and \((\partial u)\,\zeta\,n\). We estimate the two pieces.

\emph{On \(\{|x|=\delta\}\cap\operatorname{supp}\zeta\):} the surface measure is \(d\mathcal H^3=\delta\,d\vartheta\,dy\) with \(|y|<R\), of total mass \(O(\delta)\). There \(|u|\le C(R)\) and \(\int_{|y|<R}|\nabla u|\,dy\le C(R)\int_{|y|<R}|y|^{-\theta}dy<\infty\), so both kinds of terms are \(O(\delta)\).

\emph{On \(\{|y|=\delta\}\cap\operatorname{supp}\zeta\):} the surface measure is \(O(\delta)\); there \(|u|\le C(R)\) and \(|\nabla u|\le C(R)\delta^{-\theta}\), so the terms are \(O(\delta)\) and \(O(\delta^{1-\theta})\) respectively.

Since \(\theta<1\), \(R_\delta\to0\) as \(\delta\downarrow0\). Both integrands are dominated by \(L^1\) functions (\(|u|\,|\partial\partial\zeta|\) is bounded with compact support; \(|D^2u|\,|\zeta|\) is in \(L^1\) as shown), so the integrals converge to those over \(\mathbb R^4\). Thus the distributional Hessian is the pointwise one. The same argument with one integration by parts (only the terms \(u(\partial\zeta)n\) occur, of size \(O(\delta)\)) identifies the distributional gradient.
\end{proof}

\begin{lemma}[smoothed Bellman function]\label{lem:2-6}
Let \(\rho\in C_c^\infty(\mathbb R^4)\) be nonnegative, supported in the unit
ball, and normalized by \(\int_{\mathbb R^4}\rho=1\).  Set
\(\rho_\varepsilon(z)=\varepsilon^{-4}\rho(z/\varepsilon)\) and
\(u_\varepsilon=u*\rho_\varepsilon\).  Then:

\begin{enumerate}
\def\labelenumi{\arabic{enumi}.}
\item
  \(u_\varepsilon\in C^\infty(\mathbb R^4)\); for some \(C=C(M,p)\),
  \[|u_\varepsilon(z)|\le C\bigl(|z|^p+\varepsilon^p\bigr),\]
  and \(\nabla u_\varepsilon\) is bounded on compact subsets of \(\mathbb R^4\).
\item
  For every \(z\in\mathbb R^4\) and every \(a,b\in\mathbb C\), with
  \[v_1=(a,a-ib),\qquad v_2=(b,-i(a-ib)),\]
  one has
  \[\sum_{j=1}^2\big\langle D^2u_\varepsilon(z)v_j,v_j\big\rangle\le0. \tag{2.5}\]
\item
  \(u_\varepsilon\to u\) locally uniformly as \(\varepsilon\downarrow0\).
\end{enumerate}
\end{lemma}

\begin{proof}
By Lemma 2.5, \(u\in W^{2,1}_{\mathrm{loc}}(\mathbb R^4)\) and has polynomial
growth.  Standard properties of convolution therefore give
\(u_\varepsilon\in C^\infty\),
\[D^2u_\varepsilon=(D^2u)*\rho_\varepsilon,
\qquad \nabla u_\varepsilon=(\nabla u)*\rho_\varepsilon.\]
Because \(\rho_\varepsilon\) is supported in \(B(0,\varepsilon)\), the growth
bound follows from
\(|u(z-\omega)|\le C(|z|+|\omega|)^p\), and the local boundedness of the
gradient is immediate from smoothness.

Fix \(z,a,b\).  For almost every \(\omega\), the point \(z-\omega=(x',y')\)
satisfies \(x'\ne0\ne y'\), so Proposition 2.3 gives
\[\sum_{j=1}^2
  \big\langle D^2u(z-\omega)v_j,v_j\big\rangle\le0.\]
The vectors \(v_1,v_2\) do not depend on the base point.  Integrating the last
inequality against the nonnegative kernel \(\rho_\varepsilon\) proves (2.5).
Finally, continuity of \(u\) implies local uniform convergence under this
compactly supported approximate identity.
\end{proof}

\subsection{The criterion in the variable \(t=r/s\)}\label{subsec:profile-criterion}

It is convenient, for the explicit examples of Section~4, to record (A3) in terms of the profile of \(\varphi\) along \(s=1\). By homogeneity we may write \[\varphi(r,s)=s^{p}f(t),\qquad t:=\frac rs\in(0,\infty),\qquad f\in C^\infty(0,\infty). \tag{2.6}\]

\begin{lemma}\label{lem:2-7}
With (2.6) and \(r,s>0\), \[\varphi_r=s^{p-1}f',\qquad
\varphi_{rr}=s^{p-2}f'',\qquad
\frac{\varphi_r}{r}=s^{p-2}\frac{f'}{t},\] \[\varphi_{rs}=s^{p-2}\big[(p-1)f'-tf''\big],\qquad
\frac{\varphi_s}{s}=s^{p-2}\big[pf-tf'\big],\qquad
\varphi_{ss}=s^{p-2}\big[p(p-1)f-2(p-1)tf'+t^2f''\big],\] all evaluated at \(t=r/s\). Consequently \[\mathcal A=\frac{s^{p-2}}{t}\,\mathsf P_1,\qquad
\mathcal C=\frac{s^{p-2}}{t}\,\mathsf P_2,\qquad
\text{and }\ \mathcal B=\frac{s^{p-2}}{t}\,\mathsf P_3\ \text{ provided }\ \mathsf Q\ge0,\] where \[\mathsf P_1:=tf''+f',\qquad
\mathsf P_2:=p^2tf-(2p-1)t^2f'+t^3f''+f',\qquad
\mathsf P_3:=t^2f''-(p-1)tf'-f',\] \[\mathsf Q:=t^2f''-(p-1)tf' .\] Hence, if for all \(t>0\) \[f'\le0,\qquad \mathsf Q\ge0,\qquad \mathsf P_1\le0,\qquad \mathsf P_2\le0,\qquad
\mathsf R:=\mathsf P_1\mathsf P_2-\mathsf P_3^2\ \ge\ 0, \tag{2.7}\] then (D0) and (D1) hold at every \((r,s)\) with \(r,s>0\).
\end{lemma}

\begin{proof}
Differentiate (2.6), using \(\partial_r t=1/s\) and \(\partial_s t=-t/s\): \(\varphi_r=s^p f'\cdot\frac1s=s^{p-1}f'\) and \(\varphi_{rr}=s^{p-1}f''\cdot\frac1s=s^{p-2}f''\), giving also \(\varphi_r/r=s^{p-1}f'/(st)=s^{p-2}f'/t\). Next \[\varphi_s=ps^{p-1}f+s^pf'\cdot\Big(-\frac ts\Big)=s^{p-1}\big[pf-tf'\big],\] so \(\varphi_s/s=s^{p-2}[pf-tf']\). Differentiating \(\varphi_r=s^{p-1}f'\) in \(s\), \[\varphi_{rs}=(p-1)s^{p-2}f'+s^{p-1}f''\cdot\Big(-\frac ts\Big)=s^{p-2}\big[(p-1)f'-tf''\big].\] Differentiating \(\varphi_s=s^{p-1}[pf-tf']\) in \(s\) and using \(\partial_s\big(pf-tf'\big)=pf'\cdot(-\tfrac ts)+\tfrac ts f'-t f''\cdot(-\tfrac ts)
=\tfrac1s\big[-(p-1)tf'+t^2f''\big]\), \[\varphi_{ss}=(p-1)s^{p-2}\big[pf-tf'\big]+s^{p-2}\big[-(p-1)tf'+t^2f''\big]
=s^{p-2}\big[p(p-1)f-2(p-1)tf'+t^2f''\big].\] Now \[\mathcal A=\varphi_{rr}+\frac{\varphi_r}{r}=s^{p-2}\Big[f''+\frac{f'}{t}\Big]
=\frac{s^{p-2}}{t}\big[tf''+f'\big]=\frac{s^{p-2}}{t}\mathsf P_1,\] \[\mathcal C=\varphi_{ss}+\frac{\varphi_s}{s}+\frac{\varphi_r}{r}
=s^{p-2}\Big[p(p-1)f-2(p-1)tf'+t^2f''+pf-tf'+\frac{f'}{t}\Big]
=\frac{s^{p-2}}{t}\big[p^2tf-(2p-1)t^2f'+t^3f''+f'\big],\] which is \(\frac{s^{p-2}}{t}\mathsf P_2\). Finally \(\mathsf Q\ge0\) means \(t^2f''-(p-1)tf'\ge0\), i.e.~(dividing by \(t>0\)) \(tf''-(p-1)f'\ge0\), i.e. \(\varphi_{rs}=s^{p-2}[(p-1)f'-tf'']\le0\); in that case \[\mathcal B=|\varphi_{rs}|-\frac{\varphi_r}{r}=s^{p-2}\Big[tf''-(p-1)f'-\frac{f'}{t}\Big]
=\frac{s^{p-2}}{t}\big[t^2f''-(p-1)tf'-f'\big]=\frac{s^{p-2}}{t}\mathsf P_3 .\] Since \(s^{p-2}/t>0\), the conditions \(\mathcal A\le0\), \(\mathcal C\le0\), \(\mathcal B^2\le\mathcal A\mathcal C\) are equivalent to \(\mathsf P_1\le0\), \(\mathsf P_2\le0\), \(\mathsf P_3^2\le\mathsf P_1\mathsf P_2\), and (D0) is \(f'\le0\).
\end{proof}

\subsection{Generalized polynomial profiles}\label{subsec:polynomial-profiles}

\begin{lemma}\label{lem:2-8}
Let \(p\ge2\), let \(E\subset[2,p)\) be a finite set of \textbf{real} exponents, let \(b_k\ge0\) for \(k\in E\), let \(a_0>0\), and put \[f(t)=a_0-\sum_{k\in E} b_kt^k\qquad(t>0),\qquad
\varphi(r,s):=s^pf(r/s)=a_0s^p-\sum_{k\in E}b_k\,r^ks^{p-k} .\] Then, with \(\mathsf P_i,\mathsf Q\) as in Lemma 2.7, \[\mathsf P_1=-\sum_{k}k^2b_k\,t^{k-1}\ \le\ 0,\qquad
\mathsf P_3=\sum_{k}k\,b_k\,t^{k-1}+\sum_{k}k(p-k)\,b_k\,t^{k}\ \ge\ 0,\] \[\mathsf P_2=p^2a_0\,t-\sum_{k}k\,b_k\,t^{k-1}-\sum_{k}(p-k)^2b_k\,t^{k+1},\] \[\mathsf Q=\sum_{k}k(p-k)\,b_k\,t^{k}\ \ge\ 0 ,\qquad f'\le0 .\] Moreover \(\varphi\) extends continuously to \([0,\infty)^2\) with \(\varphi(r,0)=0\) (so (A5) holds), is positively homogeneous of degree \(p\), is \(C^\infty\) on \(\{r>0,s>0\}\) --- i.e.~(A1) holds --- and satisfies (A2) with \[\theta:=\max\Big\{0,\ \max_{k\in E}(k+1-p)\Big\}\ \in\ [0,1).\]
\end{lemma}

\begin{proof}
\(f'=-\sum k b_kt^{k-1}\le0\) and \(f''=-\sum k(k-1)b_kt^{k-2}\). Hence \[\mathsf P_1=tf''+f'=-\sum k(k-1)b_kt^{k-1}-\sum kb_kt^{k-1}=-\sum k^2b_kt^{k-1},\] \[\mathsf P_3=t^2f''-(p-1)tf'-f'=-\sum k(k-1)b_kt^{k}+(p-1)\sum kb_kt^{k}+\sum kb_kt^{k-1},\] and \(-(k-1)+(p-1)=p-k\) gives the stated form; likewise \(\mathsf Q=t^2f''-(p-1)tf'=\sum k\big[(p-1)-(k-1)\big]b_kt^k=\sum k(p-k)b_kt^k\), which is \(\ge0\) because \(k<p\) for every \(k\in E\) (this is also why \(\mathsf P_3\ge0\)). For \(\mathsf P_2\), \[\mathsf P_2=p^2t\Big(a_0-\sum b_kt^k\Big)+(2p-1)\sum kb_kt^{k+1}-\sum k(k-1)b_kt^{k+1}-\sum kb_kt^{k-1},\] and \(-p^2+(2p-1)k-k(k-1)=-(p^2-2pk+k^2)=-(p-k)^2\), which gives the stated form.

\emph{Continuity and (A5).} \(\varphi(r,s)=a_0s^p-\sum b_kr^ks^{p-k}\) and every exponent \(p-k>0\), so each term tends to \(0\) as \(s\downarrow0\) with \(r\) fixed; thus \(\varphi\) extends continuously by \(\varphi(r,0)=0\) and \(\varphi(0,s)=a_0s^p\). Homogeneity and smoothness on the open quadrant are clear (for \(r>0\), \(r\mapsto r^k\) is \(C^\infty\) for any real \(k\)).

\emph{(A2).} Each of \(\varphi\), \(\varphi_r,\varphi_s\), \(\varphi_{rr},\varphi_{rs},\varphi_{ss},\varphi_r/r,\varphi_s/s\) is a finite linear combination, with coefficients depending only on \(p\), \(a_0\) and \(\max_kb_k\), of monomials \(r^\alpha s^\beta\) with \(\alpha\ge0\) and \(\alpha+\beta=p\), \(p-1\), \(p-2\) respectively. Explicitly: \[\varphi_r=-\sum kb_kr^{k-1}s^{p-k},\qquad
\varphi_s=pa_0s^{p-1}-\sum(p-k)b_kr^ks^{p-k-1},\] \[\varphi_{rr}=-\sum k(k-1)b_kr^{k-2}s^{p-k},\qquad
\frac{\varphi_r}{r}=-\sum kb_kr^{k-2}s^{p-k},\qquad
\varphi_{rs}=-\sum k(p-k)b_kr^{k-1}s^{p-k-1},\] \[\varphi_{ss}=p(p-1)a_0s^{p-2}-\sum(p-k)(p-k-1)b_kr^ks^{p-k-2},\qquad
\frac{\varphi_s}{s}=pa_0s^{p-2}-\sum(p-k)b_kr^ks^{p-k-2}.\] All exponents \(\alpha\) of \(r\) occurring above are \(\ge0\), because \(k\ge2\) for every \(k\in E\). The exponents \(\beta\) of \(s\) occurring are, listed by order \(m\):

\begin{itemize}
\tightlist
\item
  \(m=0\): \(\beta=p\) and \(\beta=p-k>0\). Hence \(\beta\ge0\).
\item
  \(m=1\): \(\beta=p-k>0\) (in \(\varphi_r\)), \(\beta=p-1>0\) and \(\beta=p-k-1\) (in \(\varphi_s\)). Since \(\theta\ge k+1-p\) for every \(k\in E\), we get \(p-k-1\ge-\theta\); hence \(\beta\ge-\theta\).
\item
  \(m=2\): \(\beta=p-k>0\) (in \(\varphi_{rr},\varphi_r/r\)), \(\beta=p-k-1\ge-\theta\) (in \(\varphi_{rs}\)), \(\beta=p-2\ge0\) and \(\beta=p-k-2\ge-1-\theta\) (in \(\varphi_{ss},\varphi_s/s\)). Hence \(\beta\ge-(1+\theta)\).
\end{itemize}

Note \(\theta<1\) because \(k<p\) for all \(k\in E\).

Now let \(r^\alpha s^\beta\) with \(\alpha\ge0\), \(\alpha+\beta=p-m\) and \(\beta\ge-\gamma\) for some \(\gamma\ge0\). Since \(r\le\sigma\) and \(s\le\sigma\), \[r^\alpha s^\beta\le\sigma^\alpha s^\beta=
\begin{cases}\sigma^{\alpha}s^{\beta}\le\sigma^{p-m}, &\beta\ge0,\\[2pt]
\sigma^{p-m-\beta}s^{\beta}=\sigma^{p-m}\big(\tfrac\sigma s\big)^{-\beta}
\le\sigma^{p-m}\big(\tfrac\sigma s\big)^{\gamma}, &\beta<0,\end{cases}\] using \(\sigma/s\ge1\). Applying this with \((m,\gamma)=(0,0)\), \((1,\theta)\) and \((2,1+\theta)\) gives exactly the three bounds of (A2), with \(M\) a constant depending only on \(p\), \(a_0\) and \(\{b_k\}\).
\end{proof}

\begin{remark}\label{rem:2-9}
Lemma 2.8 permits exponents \(k\) in the whole range \(2\le k<p\); when \(k>p-1\) the derivative \(\varphi_s\) blows up like \(s^{-(k+1-p)}\) as \(s\downarrow0\), which is precisely what the parameter \(\theta\) in (A2) accommodates. This extra freedom is used at \(p=\frac{16}3\) in Section~4.2 (there \(E=\{2,3,4,5\}\) and \(5>p-1=\frac{13}3\), so \(\theta=\frac23\)); at \(p=\frac{15}2\) in Section~4.3 all exponents satisfy \(k\le p-1\) and \(\theta=0\).
\end{remark}

\begin{remark}[The quadratic coefficient]\label{rem:2-10} Assume \(b_2>0\).  As \(t\downarrow0\),
\[\mathsf P_1=-4b_2t+o(t),\qquad
  \mathsf P_2=(p^2a_0-2b_2)t+o(t),\qquad
  \mathsf P_3=2b_2t+o(t),\]
and hence
\[\mathsf R=4b_2(b_2-p^2a_0)t^2+o(t^2).\]
Thus \(\mathsf R\ge0\) near the origin forces
\(b_2\ge p^2a_0\).  Both explicit profiles below saturate this necessary
condition, \(b_2=p^2a_0\), which makes the leading term vanish.  This local
observation is not, by itself, a global optimality statement.
\end{remark}

\begin{lemma}[majorization by weighted AM--GM]\label{lem:2-11}
In the setting of Lemma 2.8, let \(t_0>0\), \(\lambda>0\) and \(\mathsf M>0\) satisfy the two inequalities \[\lambda\Big(a_0-\sum_{k\in E}b_k\Big(1-\frac kp\Big)t_0^{\,k}\Big)\ \ge\ 1,
\tag{M1}\] \[\mathsf M\ \ge\ \frac{\lambda}{p}\sum_{k\in E}k\,b_k\,t_0^{\,k-p}. \tag{M2}\] Then \(\lambda\varphi\) satisfies (A4) with \(\Lambda:=\mathsf M^{1/p}\), i.e. \(\lambda\varphi(r,s)\ge s^p-\Lambda^pr^p\) for all \(r,s\ge0\).
\end{lemma}

\begin{proof}
Step 1 (Young's inequality). For \(0<\vartheta\le1\) and \(X\ge0\), \[X^{\vartheta}\ \le\ \vartheta X+1-\vartheta .\] Indeed \(h(X):=\vartheta X+1-\vartheta-X^\vartheta\) satisfies \(h(1)=0\) and \(h'(X)=\vartheta(1-X^{\vartheta-1})\), which is \(\le0\) for \(0<X\le1\) and \(\ge0\) for \(X\ge1\) (because \(\vartheta-1\le0\)); so \(h\) attains its minimum on \([0,\infty)\) at \(X=1\), where it vanishes. (At \(X=0\): \(h(0)=1-\vartheta\ge0\).)

Step 2. Let \(k\in E\), so \(0<k<p\). Apply Step 1 with \(\vartheta=k/p\in(0,1)\) and \(X=(t/t_0)^p\), noting \(X^\vartheta=(t/t_0)^k\): \[\Big(\frac t{t_0}\Big)^{k}\ \le\ \frac kp\Big(\frac t{t_0}\Big)^{p}+1-\frac kp
\qquad\Longrightarrow\qquad
t^{k}\ \le\ \frac kp\,t_0^{\,k-p}\,t^{p}+\Big(1-\frac kp\Big)t_0^{\,k}
\qquad(t\ge0). \tag{2.8}\]

Step 3. Set \(G(t):=\mathsf Mt^p+\lambda f(t)-1\). By (2.8) and \(b_k\ge0\), \[\lambda\sum_kb_kt^k\ \le\
\Big(\frac{\lambda}{p}\sum_kk\,b_k\,t_0^{\,k-p}\Big)t^p
+\lambda\sum_kb_k\Big(1-\frac kp\Big)t_0^{\,k},\] hence for all \(t\ge0\) \[G(t)=\mathsf Mt^p+\lambda a_0-\lambda\sum_kb_kt^k-1
\ \ge\ \Big[\mathsf M-\frac{\lambda}{p}\sum_kk\,b_kt_0^{\,k-p}\Big]t^p
+\Big[\lambda\Big(a_0-\sum_kb_k\Big(1-\tfrac kp\Big)t_0^{\,k}\Big)-1\Big]\ \ge\ 0\] by (M2) and (M1).

Step 4. For \(s>0\), dividing \(\lambda\varphi(r,s)\ge s^p-\Lambda^pr^p\) by \(s^p\) and putting \(t=r/s\) turns it into \(\lambda f(t)\ge1-\mathsf Mt^p\), i.e.~\(G(t)\ge0\). For \(s=0\) it reads \(0=\lambda\varphi(r,0)\ge-\Lambda^pr^p\), which is true.
\end{proof}

Finally note that if \(\varphi\) satisfies (A1)--(A3) then so does \(\lambda\varphi\) for any \(\lambda>0\) (all of (D0), (D1) are positively homogeneous in \(\varphi\): \(\mathcal A\), \(\mathcal B\), \(\mathcal C\) scale by \(\lambda\), so \(\mathcal B^2\le\mathcal A\mathcal C\) scales by \(\lambda^2\); and (A2) holds with \(\lambda M\) and the same \(\theta\)).


\section{The probabilistic model and the Bellman principle}\label{sec:probabilistic}

Fix \(p\ge2\), \(f,g\in C_c^\infty(\mathbb R^2;\mathbb C)\) and \(T>0\); write \(U=e^{t\Delta}f\), \(V=e^{t\Delta}g\) as in Section~0. Recall that \(U,V\) and all their space derivatives are bounded and continuous on \(\mathbb R^2\times[0,T]\), with \(\|\partial^\alpha U(\cdot,t)\|_q\le\|\partial^\alpha f\|_q\).

\subsection{Brownian motion with generator \(\Delta\) and the \(\sigma\)-finite measure \(\nu\)}\label{subsec:wiener}

Let \(\Omega=C([0,T];\mathbb R^2)\) with the coordinate process \(B_s(\omega)=\omega(s)\), the (augmented, right-continuous) filtration \((\mathcal F_s)_{0\le s\le T}\), and for \(x\in\mathbb R^2\) let \(P_x\) be the law of \(x+\sqrt2\,W\), \(W\) a standard planar Brownian motion. Thus under \(P_x\) \[d\langle B^i,B^j\rangle_s=2\delta_{ij}\,ds,\qquad
P_x(B_s\in dy)=G_s(y-x)\,dy , \tag{3.1}\] so the transition semigroup of \(B\) is exactly the heat semigroup \(e^{s\Delta}\) and the generator is \(\Delta\). The map \(x\mapsto E_x[\Phi]\) is Borel measurable for every bounded measurable \(\Phi\) (standard for this Feller family).

Define the \textbf{\(\sigma\)-finite} measure on \((\Omega,\mathcal F_T)\) \[\nu(F):=\int_{\mathbb R^2}P_x(F)\,dx,\qquad
E_\nu[\Phi]:=\int_{\mathbb R^2}E_x[\Phi]\,dx\] (defined for \(\Phi\ge0\) measurable, and for \(\Phi\in L^1(\nu)\)). It is \(\sigma\)-finite because \(\nu(\{B_0\in B_R\})=|B_R|<\infty\) and these sets exhaust \(\Omega\).

\begin{lemma}[stationarity]\label{lem:3-1}
For every \(s\in[0,T]\) and every Borel \(\phi\ge0\) on \(\mathbb R^2\), \(E_\nu[\phi(B_s)]=\int_{\mathbb R^2}\phi\,dx\).
\end{lemma}

\begin{proof}
By (3.1) and Tonelli, \(\int E_x[\phi(B_s)]dx=\int\!\!\int\phi(y)G_s(y-x)\,dy\,dx=\int\phi(y)\big(\int G_s(y-x)dx\big)dy=\int\phi\,dy\), since \(G_s\) is a probability density. (For \(s=0\) this is trivial.)
\end{proof}

\subsection{The two martingales}\label{subsec:martingales}

Since \(U\in C^{2,1}(\mathbb R^2\times[0,T])\) and \(\partial_tU=\Delta U\), Itô's formula applied under \(P_x\) to the real and imaginary parts of \(U\) gives, for \(s\in[0,T]\), \[X_s:=U(B_s,T-s)=U(x,T)+\int_0^s\nabla U(B_r,T-r)\cdot dB_r , \tag{3.2}\] because the drift is \(\big(-\partial_tU+\Delta U\big)dr=0\) (the Itô second-order term is \(\frac12\sum_{ij}\partial_i\partial_jU\,d\langle B^i,B^j\rangle=\Delta U\,dr\) by (3.1)). Similarly \(N_s:=V(B_s,T-s)=V(x,T)+\int_0^s\nabla V(B_r,T-r)\cdot dB_r\). Define the \textbf{transformed} martingale \[Y_s:=\int_0^s\big(A\nabla U\big)(B_r,T-r)\cdot dB_r,\qquad Y_0=0 . \tag{3.3}\] All integrands are bounded continuous, so (3.2)--(3.3) are genuine \(L^2(P_x)\) martingales for each \(x\).

We record the complex Itô isometry: if \(\mathsf a,\mathsf b:[0,T]\times\Omega\to\mathbb C^2\) are bounded adapted continuous, then by (3.1) \[E_x\Big[\Big(\int_0^T\mathsf a\cdot dB\Big)\overline{\Big(\int_0^T\mathsf b\cdot dB\Big)}\Big]
=2\,E_x\int_0^T \mathsf a\cdot\overline{\mathsf b}\,dr,\qquad
\mathsf a\cdot\overline{\mathsf b}=\sum_{j}\mathsf a_j\overline{\mathsf b_j}. \tag{3.4}\] (Expand into real and imaginary parts and use the real Itô isometry \(E[\int \mathsf a\,dB\int \mathsf b\,dB]=2E\int \mathsf a\cdot \mathsf b\), valid by (3.1).)

Combining (3.4) with Lemma 3.1 and Tonelli, and using \(|A\nabla U|^2\le4|\nabla U|^2\) from (1.1) together with (1.4), \[E_\nu|Y_T|^2=2\int_0^T\!\!\int_{\mathbb R^2}|A\nabla U|^2dx\,dt\le
8\int_0^\infty\!\!\int|\nabla U|^2dx\,dt=4\|f\|_2^2<\infty , \tag{3.5}\] likewise \(E_\nu|N_T-N_0|^2\le\|g\|_2^2<\infty\) and \(E_\nu|N_0|^2=\|V(\cdot,T)\|_2^2\le\|g\|_2^2<\infty\). So \(Y_T,\;N_0,\;N_T\in L^2(\nu)\).

\textbf{The columns.} Set, for \(r\in[0,T]\), \[\mathsf a:=\partial_1U(B_r,T-r),\quad \mathsf b:=\partial_2U(B_r,T-r),\quad
\mathsf c:=\mathsf a-i\mathsf b .\] By Lemma 1.4, \(A\nabla U=(\mathsf c,-i\mathsf c)\). Hence, identifying \(\mathbb C\simeq\mathbb R^2\), the \(\mathbb R^4\)-valued continuous semimartingale \[Z_s:=(X_s,\,Y_s)\in\mathbb R^2\times\mathbb R^2\] satisfies \(Z_s=Z_0+\int_0^s\mathbb M_r\,dB_r\) with \(Z_0=(U(x,T),0)\), where the \(j\)-th column of the real \(4\times2\) matrix \(\mathbb M_r\) is \[v_1=(\mathsf a,\ \mathsf c),\qquad v_2=(\mathsf b,\ -i\mathsf c),
\qquad \mathsf c=\mathsf a-i\mathsf b . \tag{3.6}\] \textbf{This is exactly the family appearing in Proposition 2.3 and Lemma 2.6(2).}

\subsection{The Bellman estimate}\label{subsec:bellman-estimate}

\begin{proposition}\label{prop:3-2}
Let \(p\ge2\) and let \(\varphi\) be \((p,\Lambda)\)-admissible in the sense of Definition 2.4. Then for every \(T>0\) and every \(f\in C_c^\infty\), \[\Big(E_\nu|Y_T|^p\Big)^{1/p}\ \le\ \Lambda\,\|f\|_p . \tag{3.7}\]
\end{proposition}

\begin{proof}
Let \(u(x,y)=\varphi(|x|,|y|)\), fix \(\varepsilon>0\) and let \(u_\varepsilon\) be as in Lemma 2.6. Choose \(N>\|f\|_\infty\) and set \(\tau_N:=\inf\{s\in[0,T]:|Z_s|\ge N\}\wedge T\).  Then continuity gives \(|Z_{T\wedge\tau_N}|\le N\).  Itô's formula for \(u_\varepsilon(Z_s)\) gives, \(P_x\)-a.s., \[u_\varepsilon(Z_{T\wedge\tau_N})=u_\varepsilon(Z_0)
+\underbrace{\int_0^{T\wedge\tau_N}\!\!\nabla u_\varepsilon(Z_r)\cdot dZ_r}_{=:\,\mathcal M}
+\int_0^{T\wedge\tau_N}\!\!\sum_{j=1}^2\big\langle D^2u_\varepsilon(Z_r)v_j(r),v_j(r)\big\rangle\,dr ,\] where we used \(d\langle Z^{\mathsf a},Z^{\mathsf b}\rangle_r
=2\sum_j\mathbb M_{\mathsf aj}\mathbb M_{\mathsf bj}\,dr\) (from (3.1)), so that the second-order term is \(\frac12\cdot2\sum_j\langle D^2u_\varepsilon v_j,v_j\rangle\,dr\). By (3.6) and Lemma 2.6(2), that last integral is \(\le0\). Hence, pointwise on \(\Omega\), \[u_\varepsilon(Z_{T\wedge\tau_N})\ \le\ u_\varepsilon(Z_0)+\mathcal M . \tag{3.8}\]

For \(r<\tau_N\) we have \(|Z_r|\le N\), and \(\nabla u_\varepsilon\) is continuous, hence bounded on the closed ball \(\{|z|\le N\}\) by some \(C_{N,\varepsilon}<\infty\) (Lemma 2.6(1)); while \(|\mathbb M_r|\le C'\,|\nabla U(B_r,T-r)|\le C'\|\nabla f\|_\infty\). Thus \(\mathcal M\) is a stochastic integral of a bounded adapted integrand, so \(\mathcal M\in L^2(P_x)\) and \(E_x[\mathcal M]=0\) for every \(x\).

Fix \(R>0\), put \(B_R:=\{x:|x|<R\}\), and let \(\mu_R\) denote the finite measure on \(B_R\times\Omega\) induced by the probability kernel \(x\mapsto P_x\):
\[\int_{B_R\times\Omega}F\,d\mu_R
  :=\int_{B_R}E_x[F(x,\cdot)]\,dx\qquad(F\ge0).\]
Taking \(E_x\) in (3.8) (all terms are \(P_x\)-integrable, since \(|u_\varepsilon(Z_{T\wedge\tau_N})|\le C(N^p+\varepsilon^p)\)) and then integrating over \(B_R\) gives
\[\int_{B_R}E_x\big[u_\varepsilon(Z_{T\wedge\tau_N})\big]dx\ \le\
\int_{B_R}u_\varepsilon\big(U(x,T),0\big)\,dx . \tag{3.9}\]

Let \(\varepsilon\downarrow0\) in (3.9). On the left, \(u_\varepsilon\to u\) locally uniformly (Lemma 2.6(3)) and \(|Z_{T\wedge\tau_N}|\le N\), so the integrand converges pointwise and is bounded by \(C(N^p+1)\) on the finite measure space \((B_R\times\Omega,\mu_R)\); dominated convergence applies. On the right, \(u_\varepsilon(U(x,T),0)\to u(U(x,T),0)\) pointwise with \(|u_\varepsilon(U(x,T),0)|\le C(|U(x,T)|^p+1)\in L^1(B_R)\). Hence \[\int_{B_R}E_x\big[u(Z_{T\wedge\tau_N})\big]dx\ \le\ \int_{B_R}u\big(U(x,T),0\big)dx\ \le\ 0,\] the last inequality by (A5), since \(u(x',0)=\varphi(|x'|,0)\le0\).

Now apply the majorization (A4), \(u(x',y')\ge|y'|^p-\Lambda^p|x'|^p\): \[\int_{B_R}E_x\big[|Y_{T\wedge\tau_N}|^p\big]dx\ \le\ \Lambda^p\int_{B_R}E_x\big[|X_{T\wedge\tau_N}|^p\big]dx . \tag{3.10}\] The process \(X_s=U(B_s,T-s)\) is a continuous martingale bounded by \(\|f\|_\infty\), so \(|X_s|^p\) is a bounded submartingale (Jensen) and optional stopping gives \(E_x[|X_{T\wedge\tau_N}|^p]\le E_x[|X_T|^p]=E_x[|f(B_T)|^p]\). Therefore, by Lemma 3.1, \[\int_{B_R}E_x\big[|X_{T\wedge\tau_N}|^p\big]dx\le \int_{\mathbb R^2}E_x\big[|f(B_T)|^p\big]dx=\|f\|_p^p .\] Letting \(N\to\infty\) through values larger than \(\|f\|_\infty\), continuity of each path on \([0,T]\) gives \(\tau_N=T\) for all sufficiently large \(N\), a.s.; so \(|Y_{T\wedge\tau_N}|\to|Y_T|\) a.s. Fatou's lemma in (3.10) yields \[\int_{B_R}E_x\big[|Y_T|^p\big]dx\ \le\ \Lambda^p\|f\|_p^p ,\] and letting \(R\uparrow\infty\) (monotone convergence) gives (3.7).
\end{proof}

\subsection{Proof of the Bellman principle}\label{subsec:bellman-principle-proof}

\begin{proposition}\label{prop:3-3}
Let \(p\ge2\) and let \(\varphi\) be \((p,\Lambda)\)-admissible. Then for all \(f,g\in C_c^\infty\) and every \(T>0\), \[\Big|\,2\int_0^T\!\!\int_{\mathbb R^2}(A\nabla U)\cdot\overline{\nabla V}\,dx\,dt\,\Big|
\ \le\ \Lambda\,\|f\|_p\,\|g\|_{p'} .\]
\end{proposition}

\begin{proof}
Using (3.4) under each \(P_x\), then Lemma 3.1 and Fubini (legitimate since \(\int_0^T\!\int|A\nabla U||\nabla V|\le2\int_0^T\!\int|\nabla U||\nabla V|<\infty\) by Lemma 1.5 and Cauchy--Schwarz), \[E_\nu\big[Y_T\overline{(N_T-N_0)}\big]
=2\int_{\mathbb R^2}E_x\!\int_0^T (A\nabla U)\cdot\overline{\nabla V}(B_r,T-r)\,dr\,dx
=2\int_0^T\!\!\int_{\mathbb R^2}(A\nabla U)\cdot\overline{\nabla V}\,dx\,dt .\] (The expectation on the left is absolutely convergent because \(Y_T\) and \(N_T-N_0\) are in \(L^2(\nu)\) by (3.5).) Moreover \(E_x[Y_T\overline{N_0}]=\overline{V(x,T)}\,E_x[Y_T]=0\) for each \(x\), and \(Y_T\overline{N_0}\in L^1(\nu)\) by Cauchy--Schwarz in \(L^2(\nu)\); hence \(E_\nu[Y_T\overline{N_0}]=0\) and \[2\int_0^T\!\!\int(A\nabla U)\cdot\overline{\nabla V}\,dx\,dt=E_\nu\big[Y_T\overline{N_T}\big].\] By Hölder's inequality on the measure space \((\Omega,\nu)\), Proposition 3.2, and Lemma 3.1 (which gives \(\|N_T\|_{L^{p'}(\nu)}=\big(E_\nu|g(B_T)|^{p'}\big)^{1/p'}=\|g\|_{p'}\)), \[\big|E_\nu[Y_T\overline{N_T}]\big|\le\|Y_T\|_{L^p(\nu)}\|N_T\|_{L^{p'}(\nu)}
\le\Lambda\,\|f\|_p\|g\|_{p'} . \]
\end{proof}

\begin{proof}[Proof of the Bellman principle] Let \(\varphi\) be \((p,\Lambda)\)-admissible, \(p\ge2\). Let \(T\to\infty\) in Proposition 3.3: the integrand \((A\nabla U)\cdot\overline{\nabla V}\) is dominated in absolute value by \(2|\nabla U||\nabla V|\in L^1(\mathbb R^2\times(0,\infty))\) (Lemma 1.5), so dominated convergence and (1.3) give \[\Big|\int (Bf)\bar g\,dx\Big|\ \le\ \Lambda\,\|f\|_p\|g\|_{p'}\qquad(f,g\in C_c^\infty).\] Then \(Bf\in L^2\subset L^1_{\mathrm{loc}}\) and Lemma 1.2 (with \(q=p'\), \(h=Bf\), \(M=\Lambda\|f\|_p\)) gives \(\|Bf\|_p\le\Lambda\|f\|_p\). The estimate on \(C_c^\infty\) therefore extends uniquely to all of \(L^p\), and \(\|B\|_{L^p\to L^p}\le\Lambda\).
\end{proof}


\section{Three admissible Bellman functions}\label{sec:examples}

\subsection{The rescaled Burkholder function}\label{subsec:burkholder}

Fix \(p\ge2\) and write \[q:=p-1\ge1,\qquad \kappa:=\sqrt{\frac{2p}{p-1}}=\sqrt{2+\frac2q}\in(\sqrt2,\,2],
\qquad \alpha_p:=p\Big(1-\frac1p\Big)^{p-1}=\frac{(p-1)^{p-1}}{p^{p-2}}>0 .\] (Indeed \(\kappa^2=2+2/q\le4\) because \(q\ge1\), and \(\kappa^2>2\); so \(1<\sqrt2<\kappa\le2\).) Define \[\Phi(\rho,s):=\big(s-q\rho\big)(\rho+s)^{q}\qquad(\rho,s\ge0),\qquad
\varphi^{\mathrm B}(r,s):=\alpha_p\,\Phi(\kappa r,s) . \tag{4.1}\]

\begin{lemma}\label{lem:4-1}
\(\alpha_p\ge1\) for every \(p\ge2\), with equality iff \(p=2\).
\end{lemma}

\begin{proof}
Put \(X=p-1\ge1\). Then \(\alpha_p=(X+1)\big(\tfrac{X}{X+1}\big)^{X}\), and \(\alpha_p\ge1\) is equivalent (take logarithms) to \[\psi(X):=\ln(1+X)-X\ln\!\Big(1+\frac1X\Big)\ \ge\ 0 .\] We have \(\psi(1)=\ln2-\ln2=0\) and, using \(\frac{d}{dX}\ln(1+\frac1X)=\frac1{X+1}-\frac1X=-\frac1{X(X+1)}\), \[\psi'(X)=\frac1{1+X}-\ln\Big(1+\frac1X\Big)+X\cdot\frac1{X(X+1)}
=\frac{2}{1+X}-\ln\Big(1+\frac1X\Big).\] Substituting \(v=1/X\in(0,1]\), \(\psi'(X)=\chi(v):=\frac{2v}{1+v}-\ln(1+v)\). Now \(\chi(0)=0\) and \(\chi'(v)=\frac{2}{(1+v)^2}-\frac1{1+v}=\frac{1-v}{(1+v)^2}\ge0\) for \(v\in[0,1]\), so \(\chi\ge0\) on \([0,1]\), i.e.~\(\psi'\ge0\) on \([1,\infty)\) (and \(\psi'(X)>0\) for \(X>1\)). Hence \(\psi(X)\ge\psi(1)=0\) for \(X\ge1\).
\end{proof}

\begin{lemma}[majorization]\label{lem:4-2}
For all \(\rho,s\ge0\): \(\ \alpha_p\Phi(\rho,s)\ \ge\ s^p-(p-1)^p\rho^p\).
\end{lemma}

\begin{proof}
Both sides are positively homogeneous of degree \(p\) in \((\rho,s)\). The inequality is trivial at \((0,0)\); otherwise divide by \((\rho+s)^p\) and set \(t:=\rho/(\rho+s)\in[0,1]\), so \(s/(\rho+s)=1-t\). The claim becomes \(G(t)\ge0\) on \([0,1]\), where \[G(t):=\alpha_p\big(1-t-(p-1)t\big)-(1-t)^p+(p-1)^pt^p
=\alpha_p(1-pt)-(1-t)^p+(p-1)^pt^p .\]

\emph{Case \(p=2\).} Then \(\alpha_2=1\) and \(G(t)=(1-2t)-(1-t)^2+t^2=0\).

\emph{Case \(p>2\).} First, \[G\Big(\frac1p\Big)=0-\Big(1-\frac1p\Big)^{p}+\frac{(p-1)^p}{p^p}
=-\Big(\frac{p-1}{p}\Big)^{p}+\Big(\frac{p-1}{p}\Big)^{p}=0 .\] Next \(G'(t)=-\alpha_pp+p(1-t)^{p-1}+p(p-1)^pt^{p-1}\), so, using \(\alpha_p=(p-1)^{p-1}/p^{p-2}\), \[G'\Big(\frac1p\Big)=-\alpha_pp+p\Big(\frac{p-1}{p}\Big)^{p-1}+p\frac{(p-1)^p}{p^{p-1}}
=-\alpha_pp+\alpha_p+(p-1)\alpha_p=0 .\] Finally \(G''(t)=p(p-1)\big[(p-1)^pt^{p-2}-(1-t)^{p-2}\big]\), and for \(t\in[0,1]\), \(p>2\), taking \((p-2)\)-th roots of nonnegative numbers, \[G''(t)\ge0\iff (p-1)^{p/(p-2)}t\ \ge\ 1-t\iff t\ge t_0:=\frac1{1+(p-1)^{p/(p-2)}} .\] Since \(p-1>1\) and \(\frac p{p-2}>1\) we have \((p-1)^{p/(p-2)}>p-1\), hence \(1+(p-1)^{p/(p-2)}>p\) and therefore \(t_0<1/p\). Thus \(G\) is concave on \([0,t_0]\) and convex on \([t_0,1]\). On \([t_0,1]\): \(G\) is convex with \(G'(1/p)=0\) and \(1/p\in(t_0,1)\), so \(G\ge G(1/p)=0\) there. On \([0,t_0]\): \(G\) is concave, hence \(G\ge\min\{G(0),G(t_0)\}\), where \(G(t_0)\ge0\) by the previous sentence and \(G(0)=\alpha_p-1\ge0\) by Lemma 4.1.
\end{proof}

\begin{lemma}[derivatives]\label{lem:4-3}
On \(\{\rho>0,s>0\}\) put \(\sigma:=\rho+s>0\), \(\tau:=\rho/\sigma\in(0,1)\), and \[K:=p\,q\,\sigma^{p-3}>0,\qquad m:=(q-1)\tau\ \ge0,\qquad
J:=\frac{1-q\tau}{q(1-\tau)} .\] Then \[\Phi_{\rho\rho}=-K\sigma\big[1+m\big],\quad
\Phi_{\rho s}=-K\sigma\,m,\quad
\Phi_{ss}=K\sigma\big[1-m\big],\quad
\frac{\Phi_\rho}{\rho}=-K\sigma,\quad
\frac{\Phi_s}{s}=K\sigma\,J , \tag{4.2}\] and \(\Phi_\rho=-pq\rho\sigma^{p-2}\le0\). Moreover \(J\le\frac1q\).
\end{lemma}

\begin{proof}
Direct differentiation of \(\Phi(\rho,s)=(s-q\rho)(\rho+s)^{q}\): \[\Phi_\rho=-q\sigma^{q}+q(s-q\rho)\sigma^{q-1}=q\sigma^{q-1}\big[-\sigma+s-q\rho\big]
=-p\,q\,\rho\,\sigma^{p-2},\] \[\Phi_s=\sigma^{q}+q(s-q\rho)\sigma^{q-1}=\sigma^{p-2}\big[\sigma+qs-q^2\rho\big]
=p\,\sigma^{p-2}\big[s-(p-2)\rho\big]\] (using \(1-q^2=1-(p-1)^2=-p(p-2)\) and \(\sigma=\rho+s\)), and, differentiating once more, \[\Phi_{\rho\rho}=-pq\big[\sigma^{p-2}+(p-2)\rho\sigma^{p-3}\big]=-K\big[q\rho+s\big],
\qquad \Phi_{\rho s}=-pq(p-2)\rho\,\sigma^{p-3}=-K(q-1)\rho,\] \[\Phi_{ss}=p\big[(p-2)\sigma^{p-3}(s-(p-2)\rho)+\sigma^{p-2}\big]
=p\,\sigma^{p-3}\big[(p-1)s-(p-1)(p-3)\rho\big]=K\big[s-(q-2)\rho\big]\] (using \(1-(p-2)^2=-(p-1)(p-3)\)). Also \(\frac{\Phi_\rho}{\rho}=-pq\,\sigma^{p-2}=-K\sigma\) and \(\frac{\Phi_s}{s}=p\,\sigma^{p-2}\frac{s-(p-2)\rho}{s}\). Substituting \(\rho=\tau\sigma\), \(s=(1-\tau)\sigma\) and \(p\sigma^{p-2}=K\sigma/q\): \(\Phi_{\rho\rho}=-K\sigma[q\tau+1-\tau]=-K\sigma[1+m]\); \(\Phi_{\rho s}=-K\sigma(q-1)\tau=-K\sigma m\); \(\Phi_{ss}=K\sigma[(1-\tau)-(q-2)\tau]=K\sigma[1-(q-1)\tau]=K\sigma[1-m]\); and \[\frac{\Phi_s}{s}=p\sigma^{p-2}\frac{(1-\tau)-(q-1)\tau}{1-\tau}
=\frac{K\sigma}{q}\cdot\frac{1-q\tau}{1-\tau}=K\sigma J .\] Finally \(J\le1/q\iff\frac{1-q\tau}{1-\tau}\le1\iff 1-q\tau\le1-\tau\iff(q-1)\tau\ge0\), which holds.
\end{proof}

\begin{proposition}\label{prop:4-4}
For every \(p\ge2\), \(\varphi^{\mathrm B}\) is \(\big(p,\sqrt{2p(p-1)}\big)\)-admissible. Consequently, by the Bellman principle, \[\|B\|_{L^p\to L^p}\ \le\ \sqrt{2p(p-1)}\qquad(p\ge2). \tag{4.3}\]
\end{proposition}

\begin{proof}
\textbf{(A1)} \(\varphi^{\mathrm B}(r,s)=\alpha_p(s-q\kappa r)(\kappa r+s)^{q}\) is continuous on \([0,\infty)^2\), positively homogeneous of degree \(p\), and \(C^\infty\) where \(\kappa r+s>0\), in particular on the open quadrant.

\textbf{(A5)} \(\varphi^{\mathrm B}(r,0)=\alpha_p(-q\kappa r)(\kappa r)^q\le0\).

\textbf{(A4)} By Lemma 4.2 with \(\rho=\kappa r\), \(\varphi^{\mathrm B}(r,s)\ge s^p-(p-1)^p\kappa^pr^p=s^p-\Lambda^pr^p\) with \(\Lambda=\kappa(p-1)=\sqrt{2p/(p-1)}\,(p-1)=\sqrt{2p(p-1)}\).

\textbf{(A2)} \emph{(with \(\theta=0\).)} Since \(\varphi^{\mathrm B}(r,s)=\alpha_p\Phi(\rho,s)\) with \(\rho=\kappa r\), we have \[\varphi^{\mathrm B}_{rr}=\alpha_p\kappa^2\Phi_{\rho\rho},\quad
\varphi^{\mathrm B}_{rs}=\alpha_p\kappa\Phi_{\rho s},\quad
\varphi^{\mathrm B}_{ss}=\alpha_p\Phi_{ss},\quad
\frac{\varphi^{\mathrm B}_r}{r}=\alpha_p\kappa^2\frac{\Phi_\rho}{\rho},\quad
\frac{\varphi^{\mathrm B}_s}{s}=\alpha_p\frac{\Phi_s}{s} . \tag{4.4}\] Write \(\sigma'=\kappa r+s\le 2(r+s)=2\sigma\). By (4.2) and \(0\le m\le q\), \(\kappa\le2\): \(|\varphi^{\mathrm B}_{rr}|,|\varphi^{\mathrm B}_{rs}|,|\varphi^{\mathrm B}_{ss}|,
|\varphi^{\mathrm B}_r/r|\le 4\alpha_p pq(1+q)\sigma'^{\,p-2}\), while \[\Big|\frac{\varphi^{\mathrm B}_s}{s}\Big|=\alpha_p\,p\,\sigma'^{\,p-2}\frac{|s-(p-2)\kappa r|}{s}
\le\alpha_p p\big(1+2(p-2)\big)\sigma'^{\,p-2}\Big(1+\frac rs\Big).\] Also \(|\varphi^{\mathrm B}|\le\alpha_pq\sigma'^{\,p}\) (as \(|s-q\kappa r|\le q\sigma'\) for \(q\ge1\)), \(|\varphi^{\mathrm B}_r|=\alpha_p pq\kappa^2r\sigma'^{\,p-2}\le4\alpha_ppq\sigma'^{\,p-1}\) and \(|\varphi^{\mathrm B}_s|\le\alpha_pp(1+2(p-2))\sigma'^{\,p-1}\). Since \(\sigma'\le2\sigma\) and \(1+r/s=\sigma/s\), (A2) holds with \(\theta=0\) and a suitable \(M=M(p)\).

\textbf{(A3)} By (4.2) and (4.4), \(\varphi^{\mathrm B}_r=\alpha_p\kappa\Phi_\rho\le0\), which is (D0). Next, since \(m\ge0\), \(\varphi^{\mathrm B}_{rs}=-\alpha_pK\sigma'\kappa m\le0\), so \(|\varphi^{\mathrm B}_{rs}|=\alpha_pK\sigma'\kappa m\), and the entries of the drift matrix are (\(\alpha_pK\sigma'>0\) factored out): \[\frac{\mathcal A}{\alpha_pK\sigma'}=-\kappa^2(1+m)-\kappa^2=-\kappa^2(2+m),\qquad
\frac{\mathcal B}{\alpha_pK\sigma'}=\kappa m+\kappa^2=\kappa(m+\kappa),\] \[\frac{\mathcal C}{\alpha_pK\sigma'}=(1-m)+J-\kappa^2 .\] Clearly \(\mathcal A<0\). Also \(\mathcal C\le0\), because by \(J\le1/q\), \(m\ge0\) and \(\kappa^2=2+\frac2q\), \[1-m+J-\kappa^2\ \le\ 1+\frac1q-2-\frac2q=-1-\frac1q\ <\ 0 .\] It remains to check \(\mathcal B^2\le\mathcal A\mathcal C\), i.e. \(\kappa^2(m+\kappa)^2\le\kappa^2(2+m)\big(\kappa^2-1+m-J\big)\), i.e. \[D:=(m+2)\big(\kappa^2-1+m-J\big)-(m+\kappa)^2\ \ge\ 0 .\] Expanding, \((m+2)(\kappa^2-1+m)=m\kappa^2+m^2+m+2\kappa^2-2\) and \((m+\kappa)^2=m^2+2m\kappa+\kappa^2\), so \[D=m\kappa^2-2m\kappa+m+\kappa^2-2-(m+2)J
=m(\kappa-1)^2+\big(\kappa^2-2\big)-(m+2)J .\] Insert \(\kappa^2-2=\frac2q\), \(m=(q-1)\tau\), \(J=\frac{1-q\tau}{q(1-\tau)}\) and multiply by \(q(1-\tau)>0\): \[q(1-\tau)D=q(1-\tau)(q-1)\tau(\kappa-1)^2+2(1-\tau)-\big((q-1)\tau+2\big)(1-q\tau).\] The last two terms equal \[2-2\tau-\Big[(q-1)\tau-q(q-1)\tau^2+2-2q\tau\Big]
=\tau\big[-2-(q-1)+2q\big]+q(q-1)\tau^2=(q-1)\tau\big[1+q\tau\big].\] Hence \[q(1-\tau)\,D=(q-1)\,\tau\,\Big[\,q(1-\tau)(\kappa-1)^2+1+q\tau\,\Big]\ \ge\ 0 ,
\tag{4.5}\] because \(q\ge1\), \(\tau\in(0,1)\) and \((\kappa-1)^2\ge0\). Thus \(D\ge0\), (D1) holds, and \(\varphi^{\mathrm B}\) is admissible.
\end{proof}

\begin{remark}[criticality of \(\kappa\)]\label{rem:4-5}
If \(\kappa\) is replaced by \(\kappa'\in(1,\kappa)\) in (4.1) then, repeating the computation, at \(\tau\to0\) (i.e.~\(m\to0\), \(J\to1/q\)) one gets \(D\to(\kappa'^2-2)-\frac2q<0\). So \(\kappa\) cannot be lowered inside this one-parameter family; a genuinely different profile is needed, which is what Section~4.2--Section~4.3 provide.
\end{remark}

\subsection{A quintic profile at \(p=\frac{16}{3}\)}\label{subsec:quintic}

Let \(p=\frac{16}3\), so \(p-1=\frac{13}3\), \(2p-1=\frac{29}3\), \(p^2=\frac{256}9\), and put \[f(t):=729-20736\,t^2-61440\,t^3-56320\,t^4-4320\,t^5, \tag{4.6}\] \[\varphi(r,s):=s^{16/3}f(r/s)=729\,s^{16/3}-20736\,r^2s^{10/3}-61440\,r^3s^{7/3}
-56320\,r^4s^{4/3}-4320\,r^5s^{1/3}.\] This is Lemma 2.8 with \(E=\{2,3,4,5\}\subset[2,\frac{16}3)\), \(a_0=729\) and \((b_2,b_3,b_4,b_5)=(20736,\,61440,\,56320,\,4320)\); note \(b_2=p^2a_0=\frac{256}9\cdot729=256\cdot81=20736\), which saturates the necessary condition in Remark 2.10. The growth parameter is \(\theta=\max\{0,\,5+1-\tfrac{16}3\}=\tfrac23<1\). Thus (A1), (A2), (A5) hold by Lemma 2.8.

\begin{lemma}[drift]\label{lem:4-6}
\(\varphi\) satisfies (D0) and (D1) at every \((r,s)\) with \(r,s>0\); i.e.~(A3) holds.
\end{lemma}

\begin{proof}
We use Lemma 2.8 with \(p=\frac{16}3\). Since all \(b_k\ge0\) and all \(k<p\), we already know \(f'\le0\), \(\mathsf Q\ge0\), \(\mathsf P_1\le0\). Explicitly, \[\mathsf P_1=-\big(4\cdot20736\,t+9\cdot61440\,t^2+16\cdot56320\,t^3+25\cdot4320\,t^4\big)
=-8t\,\Pi_1,\] \[\Pi_1:=13500t^3+112640t^2+69120t+10368 ,\] (check: \(4\cdot20736=82944=8\cdot10368\), \(9\cdot61440=552960=8\cdot69120\), \(16\cdot56320=901120=8\cdot112640\), \(25\cdot4320=108000=8\cdot13500\)).

Next, with \(p-k=\frac{10}3,\frac73,\frac43,\frac13\) for \(k=2,3,4,5\), \[\mathsf Q=\sum_{k}k(p-k)b_kt^k
=138240\,t^2+430080\,t^3+\tfrac{901120}3t^4+7200\,t^5
=\tfrac{20}3t^2\big(1080t^3+45056t^2+64512t+20736\big)\ \ge0\] (check: \(\tfrac{20}3\cdot20736=138240\), \(\tfrac{20}3\cdot64512=430080\), \(\tfrac{20}3\cdot45056=\tfrac{901120}3\), \(\tfrac{20}3\cdot1080=7200\)).

Also \(\sum_kkb_kt^{k-1}=41472t+184320t^2+225280t^3+21600t^4\), so \[\mathsf P_3=\sum_kkb_kt^{k-1}+\mathsf Q
=41472t+322560t^2+655360t^3+\tfrac{965920}3t^4+7200\,t^5=\tfrac83t\,\Pi_3,\] \[\Pi_3:=2700t^4+120740t^3+245760t^2+120960t+15552\] (check: \(\tfrac83\cdot15552=41472\), \(\tfrac83\cdot120960=322560\), \(\tfrac83\cdot245760=655360\), \(\tfrac83\cdot120740=\tfrac{965920}3\), \(\tfrac83\cdot2700=7200\); and \(184320+138240=322560\), \(225280+430080=655360\), \(21600+\tfrac{901120}3=\tfrac{64800+901120}3=\tfrac{965920}3\)).

For \(\mathsf P_2\): \(p^2a_0=\frac{256}9\cdot729=20736\) and \((p-k)^2b_k=\frac{100}9\cdot20736=230400\), \(\frac{49}9\cdot61440=\frac{1003520}3\), \(\frac{16}9\cdot56320=\frac{901120}9\), \(\frac19\cdot4320=480\) for \(k=2,3,4,5\). Hence \[\mathsf P_2=20736t-\big(41472t+184320t^2+225280t^3+21600t^4\big)
-\big(230400t^3+\tfrac{1003520}3t^4+\tfrac{901120}9t^5+480t^6\big)\] \[=-20736t-184320t^2-455680t^3-\tfrac{1068320}3t^4-\tfrac{901120}9t^5-480t^6
=-\tfrac89t\,\Pi_2\ \le\ 0,\] \[\Pi_2:=540t^5+112640t^4+400620t^3+512640t^2+207360t+23328\] (check: \(\tfrac89\cdot23328=20736\), \(\tfrac89\cdot207360=184320\), \(\tfrac89\cdot512640=455680\), \(\tfrac89\cdot400620=\tfrac{3204960}9=\tfrac{1068320}3\), \(\tfrac89\cdot112640=\tfrac{901120}9\), \(\tfrac89\cdot540=480\); and \(225280+230400=455680\), \(21600+\tfrac{1003520}3=\tfrac{64800+1003520}3=\tfrac{1068320}3\)).

All the polynomials \(\Pi_1,\Pi_2,\Pi_3\) have nonnegative coefficients, so indeed \(\mathsf P_1\le0\), \(\mathsf P_2\le0\) for \(t>0\).

Finally, since \(\mathsf P_1\mathsf P_2=8t\cdot\frac89t\,\Pi_1\Pi_2
=\frac{64t^2}9\Pi_1\Pi_2\) and \(\mathsf P_3^2=\frac{64t^2}9\Pi_3^2\), \[\mathsf R=\mathsf P_1\mathsf P_2-\mathsf P_3^2=\frac{64t^2}9
\big(\Pi_1\Pi_2-\Pi_3^2\big).\] Expanding the two degree-\(8\) polynomials: \[\Pi_1\Pi_2=7290000t^8+1581465600t^7+18133464400t^6+59837752320t^5\] \[\qquad\qquad+89401835520t^4+63259263360t^3+22275440640t^2+3762339840t+241864704,\] \[\Pi_3^2=7290000t^8+651996000t^7+15905251600t^6+59999308800t^5\] \[\qquad\qquad+89691379200t^4+63209756160t^3+22275440640t^2+3762339840t+241864704,\] so, subtracting coefficient by coefficient (the four lowest coefficients and the leading one cancel), \[\Pi_1\Pi_2-\Pi_3^2
=2\,t^3\big(464734800t^4+1114106400t^3-80778240t^2-144771840t+24753600\big)\] and therefore \[\mathsf R=\frac{128}9\,t^{5}\,\mathsf N(t),\qquad
\mathsf N(t):=464734800t^4+1114106400t^3-80778240t^2-144771840t+24753600 . \tag{4.7}\]

It remains to prove \(\mathsf N\ge0\) on \([0,\infty)\). First we peel off, by an identity in \(t\), a part all of whose coefficients are \textbf{nonnegative} (hence \(\ge0\) for \(t\ge0\)): \[\mathsf N(t)=72000\,\mathsf M(t)+\big(46800t^4+50400t^3+5760t^2+20160t+57600\big),\] \[\mathsf M(t):=6454t^4+15473t^3-1122t^2-2011t+343 \tag{4.8}\] (check, coefficient by coefficient: \(72000\cdot6454=464688000\) and \(464688000+46800=464734800\); \(72000\cdot15473=1114056000\) and \(1114056000+50400=1114106400\); \(72000\cdot(-1122)=-80784000\) and \(-80784000+5760=-80778240\); \(72000\cdot(-2011)=-144792000\) and \(-144792000+20160=-144771840\); \(72000\cdot343=24696000\) and \(24696000+57600=24753600\)).

Now \(\mathsf M\) splits, again identically in \(t\), into a part with nonnegative coefficients plus \(\mathsf q_1+t\,\mathsf q_2\) with two nonnegative quadratics: \[\mathsf M(t)=6454t^4+44t^3+\mathsf q_1(t)+t\,\mathsf q_2(t),\qquad
\mathsf q_1(t):=5680t^2-2782t+343,\quad \mathsf q_2(t):=15429t^2-6802t+771
\tag{4.9}\] (check: the \(t^3\) coefficient is \(44+15429=15473\), the \(t^2\) coefficient is \(5680-6802=-1122\), the \(t^1\) coefficient is \(-2782+771=-2011\), and the constant term is \(343\)). Both quadratics are nonnegative \textbf{on all of \(\mathbb R\)}, by their discriminants: \[2782^2=7739524\ \le\ 4\cdot5680\cdot343=7792960,\qquad
6802^2=46267204\ \le\ 4\cdot15429\cdot771=47583036 .\] Hence \(\mathsf q_1\ge0\) and \(\mathsf q_2\ge0\) everywhere, so by (4.9) \(\mathsf M(t)\ge0\) for \(t\ge0\), so by the previous display \(\mathsf N(t)\ge0\) for \(t\ge0\), and finally \[\boxed{\ \mathsf R=\frac{128}9\,t^{5}\,\mathsf N(t)\ \ge\ 0\ \qquad(t>0).\ }\] By Lemma 2.7, (D0) and (D1) hold.
\end{proof}

\begin{lemma}[majorization]\label{lem:4-7}
Put \[t_0:=\Big(\frac{47}{88}\Big)^{3}=\frac{103823}{681472},\qquad
\lambda:=\frac{19}{6185},\qquad \mathsf M:=22830 .\] Then \(\lambda\varphi\) satisfies (A4) with \(\Lambda_1:=22830^{\,3/16}\).
\end{lemma}

\begin{proof}
We verify (M1) and (M2) of Lemma 2.11. Write \(y:=\frac{47}{88}\), so \(t_0=y^3\), \(t_0^{\,k}=y^{3k}\) and \(t_0^{\,k-p}=y^{3k-16}\).

\textbf{(M1).} Here \(1-\frac kp=1-\frac{3k}{16}\), giving the weights \(\frac58,\frac7{16},\frac14,\frac1{16}\) for \(k=2,3,4,5\), so \[b_2\Big(1-\tfrac2p\Big)=20736\cdot\tfrac58=12960,\quad
b_3\cdot\tfrac7{16}=61440\cdot\tfrac7{16}=26880,\quad
b_4\cdot\tfrac14=14080,\quad b_5\cdot\tfrac1{16}=270 ,\] and \[\mathsf S:=a_0-\sum_kb_k\Big(1-\frac kp\Big)t_0^{\,k}
=729-12960\,y^{6}-26880\,y^{9}-14080\,y^{12}-270\,y^{15}.\] Multiply by \(88^{15}=146973853900112596791703109632\) and use \(y^{3j}=47^{3j}/88^{3j}\): \[88^{15}\,\mathsf S=729\cdot88^{15}-12960\cdot47^{6}\cdot88^{9}
-26880\cdot47^{9}\cdot88^{6}-14080\cdot47^{12}\cdot88^{3}-270\cdot47^{15}.\] The five integers are \[107143939493182083061151566921728,\quad 44211596576200490352388677304320,\] \[13970309212700369991666122096640,\quad 1114871892671277389595594588160,\] \[3257104054721499424453142610,\] whence \(88^{15}\,\mathsf S=47843904707555223828076719789998\). Condition (M1), \(\lambda\mathsf S\ge1\) with \(\lambda=\frac{19}{6185}\), is therefore equivalent to \[19\cdot47843904707555223828076719789998\ \ge\
6185\cdot146973853900112596791703109632,\] i.e.~to \[909034189443549252733457676009962\ \ge\ 909033286372196411156683733073920 ,\] which is true. (The ratio is \(1.00000099\ldots\).)

\textbf{(M2).} Here \(\frac{\lambda}p=\frac{19}{6185}\cdot\frac3{16}=\frac{57}{98960}\) and \[\Sigma:=\sum_kk\,b_k\,t_0^{\,k-p}
=41472\,y^{-10}+184320\,y^{-7}+225280\,y^{-4}+21600\,y^{-1},\] using \(2b_2=41472\), \(3b_3=184320\), \(4b_4=225280\), \(5b_5=21600\). Since \(y^{-1}=\frac{88}{47}\), multiplying by \(47^{10}=52599132235830049\) gives \[
\begin{aligned}
47^{10}\Sigma
 &=41472\cdot88^{10}+184320\cdot47^{3}\cdot88^{7}\\
 &\quad+225280\cdot47^{6}\cdot88^{4}+21600\cdot47^{9}\cdot88.
\end{aligned}
\] whose four summands are
\[1154999247706192481353728,\quad 782068404177362611077120,\]
\[145626520759631683256320,\quad 2127243203273739513600.\]
Their sum is
\[47^{10}\Sigma=2084821415846460515200768.\]
Condition (M2), \(\mathsf M\ge\frac{57}{98960}\Sigma\), is therefore equivalent to
\[22830\cdot98960\cdot47^{10}\ge57\cdot2084821415846460515200768.\]
The two sides are, respectively,
\[
\begin{aligned}
22830\cdot98960\cdot47^{10}
 &=118834947177898241847583200,\\
57\cdot2084821415846460515200768
 &=118834820703248249366443776.
\end{aligned}
\]
Thus (M2) holds.  Numerically, \(\frac{\lambda}{p}\Sigma=22829.9757\ldots\le22830\).

By Lemma 2.11, \(\lambda\varphi\) satisfies (A4) with \(\Lambda_1=\mathsf M^{1/p}=22830^{3/16}\).
\end{proof}

\begin{proposition}\label{prop:4-8}
\[\|B\|_{L^{16/3}\to L^{16/3}}\ \le\ \Lambda_1=22830^{\,3/16} . \tag{4.10}\]
\end{proposition}

\begin{proof}
\(\varphi\) satisfies (A1), (A2) (with \(\theta=\frac23\)) and (A5) by Lemma 2.8, and (A3) by Lemma 4.6. Multiplying by \(\lambda=\frac{19}{6185}>0\) preserves (A1), (A2), (A3), (A5), and Lemma 4.7 gives (A4) with \(\Lambda_1\). So \(\lambda\varphi\) is \((\frac{16}3,\Lambda_1)\)-admissible, and the Bellman principle applies.
\end{proof}

Numerically \(\Lambda_1=6.5647752\ldots\), whereas \(\sqrt{2p(p-1)}=\sqrt{2\cdot\frac{16}3\cdot\frac{13}3}=\sqrt{\frac{416}9}=6.7986\ldots\); so (4.10) improves (4.3) at \(p=\frac{16}3\).

\subsection{A sextic profile at \(p=\frac{15}{2}\)}\label{subsec:sextic}

Let \(p=\frac{15}2\), so \(p-1=\frac{13}2\), \(2p-1=14\), \(p^2=\frac{225}4\), and put \[f(t):=12-675\,t^2-3300\,t^3-6600\,t^4-5423\,t^5-2640\,t^6, \tag{4.11}\] \[\varphi(r,s):=s^{15/2}f(r/s)=12\,s^{15/2}-675\,r^2s^{11/2}-3300\,r^3s^{9/2}
-6600\,r^4s^{7/2}-5423\,r^5s^{5/2}-2640\,r^6s^{3/2}.\] This is Lemma 2.8 with \(E=\{2,3,4,5,6\}\subset[2,\frac{15}2)\), \(a_0=12\) and \((b_2,\dots,b_6)=(675,3300,6600,5423,2640)\); again \(b_2=p^2a_0=\frac{225}4\cdot12=675\). All exponents satisfy \(k\le6\le p-1=\frac{13}2\), so here \(\theta=0\). Thus (A1), (A2), (A5) hold by Lemma 2.8.

\begin{lemma}[drift]\label{lem:4-9}
\(\varphi\) satisfies (D0) and (D1) at every \((r,s)\) with \(r,s>0\).
\end{lemma}

\begin{proof}
By Lemma 2.8 with \(p=\frac{15}2\): \[\mathsf P_1=-\big(4\cdot675\,t+9\cdot3300\,t^2+16\cdot6600\,t^3+25\cdot5423\,t^4
+36\cdot2640\,t^5\big)=-5t\,\Theta_1,\] \[\Theta_1:=19008t^4+27115t^3+21120t^2+5940t+540\] (check: \(5\cdot540=2700=4\cdot675\); \(5\cdot5940=29700=9\cdot3300\); \(5\cdot21120=105600=16\cdot6600\); \(5\cdot27115=135575=25\cdot5423\); \(5\cdot19008=95040=36\cdot2640\)).

With \(p-k=\frac{11}2,\frac92,\frac72,\frac52,\frac32\) for \(k=2,\dots,6\), \[
\begin{aligned}
\mathsf Q
 &=\sum_k k(p-k)b_kt^k\\
 &=7425t^2+44550t^3+92400t^4+\tfrac{135575}{2}t^5+23760t^6\\
 &=\tfrac{55}{2}t^2\bigl(864t^4+2465t^3+3360t^2+1620t+270\bigr)\ge0.
\end{aligned}
\] (check: \(\tfrac{55}2\cdot270=7425\), \(\tfrac{55}2\cdot1620=44550\), \(\tfrac{55}2\cdot3360=92400\), \(\tfrac{55}2\cdot2465=\tfrac{135575}2\), \(\tfrac{55}2\cdot864=23760\)).

Also \(\sum_kkb_kt^{k-1}=1350t+9900t^2+26400t^3+27115t^4+15840t^5\), hence \[\mathsf P_3=\sum_kkb_kt^{k-1}+\mathsf Q
=1350t+17325t^2+70950t^3+119515t^4+\tfrac{167255}2t^5+23760t^6=\tfrac52t\,\Theta_3,\] \[\Theta_3:=9504t^5+33451t^4+47806t^3+28380t^2+6930t+540\] (check: \(\tfrac52\cdot540=1350\), \(\tfrac52\cdot6930=17325\), \(\tfrac52\cdot28380=70950\), \(\tfrac52\cdot47806=119515\), \(\tfrac52\cdot33451=\tfrac{167255}2\), \(\tfrac52\cdot9504=23760\); and \(9900+7425=17325\), \(26400+44550=70950\), \(27115+92400=119515\), \(15840+\tfrac{135575}2=\tfrac{31680+135575}2=\tfrac{167255}2\)).

For \(\mathsf P_2\): \(p^2a_0=675\), and \((p-k)^2b_k=\frac{121}4\cdot675=\frac{81675}4\), \(\frac{81}4\cdot3300=66825\), \(\frac{49}4\cdot6600=80850\), \(\frac{25}4\cdot5423=\frac{135575}4\), \(\frac94\cdot2640=5940\) for \(k=2,\dots,6\). Hence \[\mathsf P_2=675t-\big(1350t+9900t^2+26400t^3+27115t^4+15840t^5\big)\] \[\qquad-\big(\tfrac{81675}4t^3+66825t^4+80850t^5+\tfrac{135575}4t^6+5940t^7\big)\] \[=-675t-9900t^2-\tfrac{187275}4t^3-93940t^4-96690t^5-\tfrac{135575}4t^6-5940t^7
=-\tfrac54t\,\Theta_2\ \le0,\] \[\Theta_2:=4752t^6+27115t^5+77352t^4+75152t^3+37455t^2+7920t+540\] (check: \(\tfrac54\cdot540=675\), \(\tfrac54\cdot7920=9900\), \(\tfrac54\cdot37455=\tfrac{187275}4\), \(\tfrac54\cdot75152=93940\), \(\tfrac54\cdot77352=96690\), \(\tfrac54\cdot27115=\tfrac{135575}4\), \(\tfrac54\cdot4752=5940\); and \(26400+\tfrac{81675}4=\tfrac{105600+81675}4=\tfrac{187275}4\), \(27115+66825=93940\), \(15840+80850=96690\)).

All of \(\Theta_1,\Theta_2,\Theta_3\) have nonnegative coefficients, so \(\mathsf P_1\le0\) and \(\mathsf P_2\le0\) for \(t>0\).

Since \(\mathsf P_1\mathsf P_2=5t\cdot\frac54t\,\Theta_1\Theta_2
=\frac{25t^2}4\Theta_1\Theta_2\) and \(\mathsf P_3^2=\frac{25t^2}4\Theta_3^2\), \[\mathsf R=\frac{25t^2}4\big(\Theta_1\Theta_2-\Theta_3^2\big).\] Expanding the two degree-\(10\) polynomials: \[\Theta_1\Theta_2=90326016t^{10}+644252400t^{9}+2305892281t^{8}+4126784376t^{7}
+4546994540t^{6}\] \[\qquad+3227458905t^{5}+1504237680t^{4}+444977280t^{3}+78675300t^{2}+7484400t+291600,\] \[\Theta_3^{\,2}=90326016t^{10}+635836608t^{9}+2027665849t^{8}+3737764052t^{7}
+4315817836t^{6}\] \[\qquad+3187363740t^{5}+1504142640t^{4}+444977280t^{3}+78675300t^{2}+7484400t+291600,\] so \[\Theta_1\Theta_2-\Theta_3^{\,2}
=t^4\big(8415792t^5+278226432t^4+389020324t^3+231176704t^2+40095165t+95040\big)\] and therefore \[\boxed{\ \mathsf R=\frac{25}4\,t^{6}\big(8415792t^5+278226432t^4+389020324t^3
+231176704t^2+40095165t+95040\big)\ \ge\ 0\ } \tag{4.12}\] for \(t>0\). Lemma 2.7 gives (D0), (D1).
\end{proof}

\begin{lemma}[majorization]\label{lem:4-10}
Put \[t_0:=\Big(\frac8{25}\Big)^{2}=\frac{64}{625},\qquad
\lambda:=\frac{105}{454},\qquad \mathsf M:=22865000 .\] Then \(\lambda\varphi\) satisfies (A4) with \(\Lambda_2:=22865000^{\,2/15}\).
\end{lemma}

\begin{proof}
Write \(w:=\frac8{25}\), so \(t_0=w^2\), \(t_0^{\,k}=w^{2k}\), \(t_0^{\,k-p}=w^{2k-15}\).

\textbf{(M1).} Here \(1-\frac kp=1-\frac{2k}{15}\), giving weights \(\frac{11}{15},\frac35,\frac7{15},\frac13,\frac15\) for \(k=2,\dots,6\), so \[675\cdot\tfrac{11}{15}=495,\quad 3300\cdot\tfrac35=1980,\quad
6600\cdot\tfrac7{15}=3080,\quad 5423\cdot\tfrac13=\tfrac{5423}3,\quad
2640\cdot\tfrac15=528,\] \[\mathsf S=12-495\,w^4-1980\,w^6-3080\,w^8-\tfrac{5423}3\,w^{10}-528\,w^{12}.\] Multiply by \(3\cdot25^{12}=178813934326171875\): \[3\cdot25^{12}\mathsf S=36\cdot25^{12}-1485\cdot8^{4}\cdot25^{8}-5940\cdot8^{6}\cdot25^{6}
-9240\cdot8^{8}\cdot25^{4}-5423\cdot8^{10}\cdot25^{2}-1584\cdot8^{12},\] whose six terms are \[2145767211914062500,\quad 928125000000000000,\quad 380160000000000000,\] \[60555264000000000,\quad 3639313694720000,\quad 108851651149824,\] giving \(3\cdot25^{12}\mathsf S=773178782568192676\). Condition (M1), \(\frac{105}{454}\mathsf S\ge1\), is therefore equivalent to \[105\cdot773178782568192676\ \ge\ 454\cdot178813934326171875,\] i.e.~to \[81183772169660230980\ \ge\ 81181526184082031250,\] which is true. (The ratio is \(1.0000276\ldots\).)

\textbf{(M2).} Here \(\frac\lambda p=\frac{105}{454}\cdot\frac2{15}=\frac{7}{227}\) and, using \(2b_2=1350\), \(3b_3=9900\), \(4b_4=26400\), \(5b_5=27115\), \(6b_6=15840\), \[\Sigma=\sum_kk\,b_k\,t_0^{\,k-p}
=1350\,w^{-11}+9900\,w^{-9}+26400\,w^{-7}+27115\,w^{-5}+15840\,w^{-3}.\] Since \(w^{-1}=\frac{25}8\), multiplying by \(8^{11}=8589934592\) gives \[8^{11}\Sigma=1350\cdot25^{11}+9900\cdot8^{2}\cdot25^{9}+26400\cdot8^{4}\cdot25^{7}
+27115\cdot8^{6}\cdot25^{5}+15840\cdot8^{8}\cdot25^{3},\] whose five terms are \[3218650817871093750,\quad 2416992187500000000,\quad 660000000000000000,\] \[69414400000000000,\quad 4152360960000000,\] totalling \(8^{11}\Sigma=6369209766331093750\). Condition (M2), \(\mathsf M\ge\frac7{227}\Sigma\), is therefore equivalent to \[22865000\cdot227\cdot8^{11}\ \ge\ 7\cdot6369209766331093750 .\] Now \(227\cdot8^{11}=1949915152384\) and \(22865000\cdot1949915152384=44584809959260160000\), while \(7\cdot6369209766331093750=44584468364317656250\); since \[44584809959260160000\ \ge\ 44584468364317656250 ,\] (M2) holds. (Indeed \(\frac\lambda p\Sigma=22864824.8\ldots\le22865000\).)

By Lemma 2.11, \(\lambda\varphi\) satisfies (A4) with \(\Lambda_2=\mathsf M^{1/p}=22865000^{2/15}\).
\end{proof}

\begin{proposition}\label{prop:4-11}
\[\|B\|_{L^{15/2}\to L^{15/2}}\ \le\ \Lambda_2=22865000^{\,2/15} . \tag{4.13}\]
\end{proposition}

\begin{proof}
Identical to Proposition 4.8: \(\varphi\) satisfies (A1), (A2) (with \(\theta=0\)) and (A5) by Lemma 2.8, and (A3) by Lemma 4.9; multiplying by \(\lambda=\frac{105}{454}>0\) preserves these, and Lemma 4.10 supplies (A4). Apply the Bellman principle.
\end{proof}

Numerically \(\Lambda_2=9.5768531\ldots\), whereas \(\sqrt{2\cdot\frac{15}2\cdot\frac{13}2}=\sqrt{\frac{195}2}=9.8742\ldots\).

\subsection{Endpoint summary}\label{subsec:endpoint-summary}

Collecting (0.2), Lemma 1.3, (4.3), (4.10) and (4.13):

\begin{boundsummary}
\(\|B\|_p=\|B\|_{p'}\) for all \(1<p<\infty\); \(\|B\|_2=1\); and \[\|B\|_p\le\sqrt{2p(p-1)}\ \ (p\ge2),\qquad
\|B\|_{16/3}\le 22830^{\,3/16},\qquad
\|B\|_{15/2}\le 22865000^{\,2/15}.\]
\end{boundsummary}


\section{Interpolation and the uniform constant}\label{sec:interpolation}

Throughout this section \(u:=1.523958\). We must show \[\|B\|_p\ \le\ u\,(p-1)\qquad\text{for all }p\ge2 ; \tag{5.1}\] the range \(1<p<2\) then follows from \(\|B\|_p=\|B\|_{p'}\) (Lemma 1.3) and \(p^*-1=p'-1\).

\subsection{A single interpolation interval}\label{subsec:one-range}

\begin{lemma}\label{lem:5-1}
Let \(2\le p_0<p_1<\infty\) and let \(L_0<L_1\) be real numbers with \(\|B\|_{p_0}\le e^{L_0}\) and \(\|B\|_{p_1}\le e^{L_1}\). Put \[\hat c:=\frac{L_1-L_0}{2\big(\frac1{p_0}-\frac1{p_1}\big)}>0,\qquad
K:=L_0-2\hat c\Big(\frac12-\frac1{p_0}\Big).\] Then for every \(p\in[p_0,p_1]\), writing \(t:=1-\frac2p\in\big[1-\frac2{p_0},\,1-\frac2{p_1}\big]\), \[\ln\frac{\|B\|_p}{p-1}\ \le\ K+\hat c\,t-2\operatorname{artanh}t . \tag{5.2}\]
\end{lemma}

\begin{proof}
\(B\) is a linear operator bounded on \(L^{p_0}\) and on \(L^{p_1}\) (Section~0), so Riesz--Thorin applies: if \(\frac1p=\frac{1-\vartheta}{p_0}+\frac{\vartheta}{p_1}\) with \(\vartheta\in[0,1]\), then \(\|B\|_p\le\|B\|_{p_0}^{1-\vartheta}\|B\|_{p_1}^{\vartheta}\le e^{(1-\vartheta)L_0+\vartheta L_1}\). Writing \(x:=1/p\), the exponent is the value at \(x\) of the affine function through \((1/p_0,L_0)\) and \((1/p_1,L_1)\), namely \[L(x)=L_0+m\Big(x-\frac1{p_0}\Big),\qquad
m:=\frac{L_1-L_0}{\frac1{p_1}-\frac1{p_0}}=-2\hat c<0 .\] With \(t=1-2x\), i.e.~\(x=\frac{1-t}2\), \[L(x)=L_0+m\Big(\frac12-\frac1{p_0}\Big)-\frac m2\,t=K+\hat c\,t .\] Finally \(p-1=\frac{1+t}{1-t}\), so \(\ln(p-1)=\ln\frac{1+t}{1-t}=2\operatorname{artanh}t\), and \(\ln\|B\|_p-\ln(p-1)\le K+\hat ct-2\operatorname{artanh}t\).
\end{proof}

\begin{lemma}[Legendre transform]\label{lem:5-2}
Let \(\hat c>2\) and \(S(t):=\hat ct-2\operatorname{artanh}t\) on \([0,1)\). Put \(g:=\sqrt{\hat c(\hat c-2)}>0\) and \(t^*:=g/\hat c=\sqrt{1-2/\hat c}\in(0,1)\). Then \(S\) is strictly increasing on \([0,t^*]\), strictly decreasing on \([t^*,1)\), and \[\max_{[0,1)}S=S(t^*)=g-2\operatorname{artanh}(t^*).\]
\end{lemma}

\begin{proof}
\(S'(t)=\hat c-\frac{2}{1-t^2}\) is strictly decreasing on \([0,1)\), with \(S'(0)=\hat c-2>0\) and \(S'(t)\to-\infty\) as \(t\uparrow1\). Its unique zero is at \(1-t^2=2/\hat c\), i.e.~\(t=t^*\); note \(\hat c\,t^*=\hat c\sqrt{1-2/\hat c}=\sqrt{\hat c^2-2\hat c}=g\), so \(S(t^*)=g-2\operatorname{artanh}(t^*)\).
\end{proof}

\begin{lemma}[shift of \(\operatorname{artanh}\), and the logarithm identity]\label{lem:5-3}
\begin{enumerate}
\def\labelenumi{\arabic{enumi}.}
\tightlist
\item
  For \(\alpha,\beta\in(-1,1)\), \(\operatorname{artanh}\alpha-\operatorname{artanh}\beta
  =\operatorname{artanh}\frac{\alpha-\beta}{1-\alpha\beta}\).
\item
  For all \(A,B>0\), \(\ \ln A=\ln B+2\operatorname{artanh}\dfrac{A-B}{A+B}\).
\item
  \(\operatorname{artanh}\tfrac35=\ln2\) and \(\operatorname{artanh}\tfrac23=\tfrac12\ln5\).
\end{enumerate}
\end{lemma}

\begin{proof}
(1) Both sides are \(\frac12\ln\) of \(\frac{(1+\alpha)(1-\beta)}{(1-\alpha)(1+\beta)}\), since \(\frac{1+\frac{\alpha-\beta}{1-\alpha\beta}}{1-\frac{\alpha-\beta}{1-\alpha\beta}}
=\frac{1-\alpha\beta+\alpha-\beta}{1-\alpha\beta-\alpha+\beta}
=\frac{(1+\alpha)(1-\beta)}{(1-\alpha)(1+\beta)}\). (2) Put \(x:=\frac{A-B}{A+B}\in(-1,1)\). Then \(\frac{1+x}{1-x}=\frac{(A+B)+(A-B)}{(A+B)-(A-B)}=\frac{2A}{2B}=\frac AB\), so \(2\operatorname{artanh}x=\ln\frac AB\). (3) \(\operatorname{artanh}\frac35=\frac12\ln\frac{1+3/5}{1-3/5}=\frac12\ln4=\ln2\) and \(\operatorname{artanh}\frac23=\frac12\ln\frac{1+2/3}{1-2/3}=\frac12\ln5\).
\end{proof}

Combining Lemmas 5.1--5.3: if \(\hat c>2\) and \(t^*\) lies in the \(t\)-interval of the range, then (5.1) holds on \([p_0,p_1]\) as soon as \[K+g\ \le\ \ln u+2\operatorname{artanh}(\beta)+2\operatorname{artanh}(\delta),
\qquad \delta:=\frac{t^*-\beta}{1-\beta t^*}, \tag{5.3}\] for any convenient base point \(\beta\in(0,1)\). Indeed, by Lemma 5.2 and Lemma 5.3(1), for every \(t\) in the range \[\ln\frac{\|B\|_p}{p-1}\le K+g-2\operatorname{artanh}t^*
=K+g-2\operatorname{artanh}\beta-2\operatorname{artanh}\delta ,\] and (5.3) says exactly that the right-hand side is \(\le\ln u\). We shall use \(\beta=\frac35\) (so \(2\operatorname{artanh}\beta=2\ln2\) and the right side of (5.3) begins with \(\ln(4u)\)) and \(\beta=\frac23\) (so \(2\operatorname{artanh}\beta=\ln5\) and it begins with \(\ln(5u)\)).

\subsection{The three ranges and their data}\label{subsec:three-ranges}

We use Theorem A at the two base points \(p=\frac{16}3\) and \(p=\frac{15}2\), together with \(\|B\|_2=1\): \[L_{(0)}:=0,\qquad
L_{(1)}:=\tfrac3{16}\ln 22830,\qquad
L_{(2)}:=\tfrac2{15}\ln 22865000 . \tag{5.4}\] These are legitimate by (0.2), (4.10) and (4.13): \(\|B\|_2=1=e^{L_{(0)}}\), \(\|B\|_{16/3}\le22830^{3/16}=e^{L_{(1)}}\), \(\|B\|_{15/2}\le22865000^{2/15}=e^{L_{(2)}}\).

\textbf{Reduction of the two logarithms to \(\ln2,\ln3,\ln5,\ln7\).} By Lemma 5.3(2):

\begin{itemize}
\tightlist
\item
  \(2^9\cdot3^7=512\cdot2187=1119744\), and \(22830\cdot49=1118670\), so with \(A=22830\), \(B=\frac{1119744}{49}\) we get \(A+B=\frac{1118670+1119744}{49}
  =\frac{2238414}{49}\), \(A-B=\frac{1118670-1119744}{49}=-\frac{1074}{49}\), and \(\frac{1074}{2238414}=\frac{179}{373069}\) (divide by \(6\)), whence \[\ln 22830=9\ln2+7\ln3-2\ln7-2\operatorname{artanh}\tfrac{179}{373069}. \tag{5.5}\]
\item
  \(2^5\cdot3^5\cdot5^5=32\cdot243\cdot3125=24300000\), and with \(A=22865000\), \(B=24300000\): \(A-B=-1435000\), \(A+B=47165000\), and \(\frac{1435000}{47165000}=\frac{287}{9433}\) (divide by \(5000\)), whence \[\ln 22865000=5\ln2+5\ln3+5\ln5-2\operatorname{artanh}\tfrac{287}{9433}. \tag{5.6}\]
\item
  \(u=1.523958\), \(B=\frac32\): \(A-B=0.023958\), \(A+B=3.023958\), so \(\frac{A-B}{A+B}=\frac{23958}{3023958}=\frac{3993}{503993}\) (divide by \(6\)) and \[\ln u=\ln3-\ln2+2\operatorname{artanh}\tfrac{3993}{503993},\] \[\ln(4u)=\ln2+\ln3+2\operatorname{artanh}\tfrac{3993}{503993},\qquad
  \ln(5u)=\ln3-\ln2+\ln5+2\operatorname{artanh}\tfrac{3993}{503993}. \tag{5.7}\]
\end{itemize}

\textbf{The three ranges.}

\begin{longtable}[]{@{}
  >{\raggedright\arraybackslash}p{(\columnwidth - 10\tabcolsep) * \real{0.1667}}
  >{\raggedright\arraybackslash}p{(\columnwidth - 10\tabcolsep) * \real{0.1667}}
  >{\raggedright\arraybackslash}p{(\columnwidth - 10\tabcolsep) * \real{0.1667}}
  >{\raggedright\arraybackslash}p{(\columnwidth - 10\tabcolsep) * \real{0.1667}}
  >{\raggedright\arraybackslash}p{(\columnwidth - 10\tabcolsep) * \real{0.1667}}
  >{\raggedright\arraybackslash}p{(\columnwidth - 10\tabcolsep) * \real{0.1667}}@{}}
\toprule\noalign{}
\begin{minipage}[b]{\linewidth}\raggedright
range
\end{minipage} & \begin{minipage}[b]{\linewidth}\raggedright
\(p\)-interval
\end{minipage} & \begin{minipage}[b]{\linewidth}\raggedright
\(t\)-interval
\end{minipage} & \begin{minipage}[b]{\linewidth}\raggedright
\(\hat c\)
\end{minipage} & \begin{minipage}[b]{\linewidth}\raggedright
\(K\)
\end{minipage} & \begin{minipage}[b]{\linewidth}\raggedright
base \(\beta\)
\end{minipage} \\
\midrule\noalign{}
\endhead
\bottomrule\noalign{}
\endlastfoot
I & \([2,\frac{16}3]\) & \([0,\frac58]\) & \(\frac85L_{(1)}=\frac3{10}\ln22830\) & \(0\) & \(\frac35\) \\
II & \([\frac{16}3,\frac{15}2]\) & \([\frac58,\frac{11}{15}]\) & \(\frac{120}{13}\big(L_{(2)}-L_{(1)}\big)\) & \(L_{(1)}-\frac58\hat c\) & \(\frac23\) \\
III & \([\frac{15}2,\infty)\) & --- & --- & --- & --- \\
\end{longtable}

(These are Lemma 5.1. Range I: \(\frac1{p_0}-\frac1{p_1}=\frac12-\frac3{16}=\frac5{16}\), so \(\hat c=\frac{L_{(1)}}{2\cdot\frac5{16}}=\frac85L_{(1)}
=\frac85\cdot\frac3{16}\ln22830=\frac3{10}\ln22830\), and \(K=0-2\hat c(\frac12-\frac12)=0\); the \(t\)-interval is \([1-\frac22,\,1-\frac{2\cdot3}{16}]=[0,\frac58]\). Range II: \(\frac3{16}-\frac2{15}=\frac{45-32}{240}=\frac{13}{240}\), so \(\hat c=\frac{L_{(2)}-L_{(1)}}{13/120}=\frac{120}{13}(L_{(2)}-L_{(1)})\), and \(K=L_{(1)}-2\hat c(\frac12-\frac3{16})=L_{(1)}-\frac58\hat c\); the \(t\)-interval is \([\frac58,\,1-\frac4{15}]=[\frac58,\frac{11}{15}]\).)

\subsection{Range III: \(p\ge\frac{15}{2}\)}\label{subsec:range3}

By (4.3), \(\frac{\|B\|_p}{p-1}\le\sqrt{\frac{2p}{p-1}}\), and \(\frac{2p}{p-1}=2+\frac2{p-1}\) is decreasing in \(p\), so for \(p\ge\frac{15}2\) (where \(p-1\ge\frac{13}2\)) \[\Big(\frac{\|B\|_p}{p-1}\Big)^2\ \le\ 2+\frac2{13/2}=\frac{30}{13} .\] Since \(13\,u^2=13\cdot(1.523958)^2=13\cdot2.322447985764=30.191823814932>30\), we get \(\|B\|_p\le u(p-1)\) on \([\frac{15}2,\infty)\).

\subsection{Rigorous numerical bounds}\label{subsec:numerics}

For \(0<x<1\) the series \(\operatorname{artanh}x=\sum_{k\ge0}\frac{x^{2k+1}}{2k+1}\) has positive terms, so every partial sum is a lower bound, and \[\sum_{k=0}^{n-1}\frac{x^{2k+1}}{2k+1}\ \le\ \operatorname{artanh}x\ \le\ \sum_{k=0}^{n-1}\frac{x^{2k+1}}{2k+1}
+\frac{x^{2n+1}}{(2n+1)(1-x^2)} \tag{5.8}\] (bounding the tail \(\sum_{k\ge n}\frac{x^{2k+1}}{2k+1}\le\frac1{2n+1}\sum_{k\ge n}x^{2k+1}
=\frac{x^{2n+1}}{(2n+1)(1-x^2)}\) by a geometric series). Using \[\ln2=2\operatorname{artanh}\tfrac13,\quad
\ln3=\ln2+2\operatorname{artanh}\tfrac15,\quad
\ln5=2\ln2+2\operatorname{artanh}\tfrac19,\quad
\ln7=3\ln2-2\operatorname{artanh}\tfrac1{15}\] (each identity is Lemma 5.3(2): \(\frac{1+x}{1-x}=2,\frac32,\frac54,\frac87\) for \(x=\frac13,\frac15,\frac19,\frac1{15}\), so \(\ln2=2\operatorname{artanh}\frac13\), \(\ln3=\ln2+\ln\frac32\), \(\ln5=\ln4+\ln\frac54\), \(\ln7=\ln8-\ln\frac87\)), together with (5.8) applied with \(n=40,25,20,18\) terms respectively, one obtains \[\ln2\in(0.6931471804,\ 0.6931471806),\qquad \ln3\in(1.0986122884,\ 1.0986122888),\] \[\ln5\in(1.6094379120,\ 1.6094379126),\qquad \ln7\in(1.9459101484,\ 1.9459101492). \tag{5.9}\] Also, directly from (5.8) with \(n=2,4,3\) terms respectively, \[\operatorname{artanh}\tfrac{179}{373069}\in(0.0004798040,\ 0.0004798041),\] \[\operatorname{artanh}\tfrac{287}{9433}\in(0.0304344966,\ 0.0304344967),\qquad
\operatorname{artanh}\tfrac{3993}{503993}\in(0.0079228948,\ 0.0079228949). \tag{5.10}\] (For instance \(x=\frac{179}{373069}=0.00047980400\ldots\), so \(x+\frac{x^3}3=0.00047980404\ldots\) and the \(n=2\) tail \(\frac{x^5}{5(1-x^2)}\) is \(<10^{-17}\); and for \(x=\frac{3993}{503993}=0.0079227290\ldots\), \(x+\frac{x^3}3+\frac{x^5}5=0.0079228948607\ldots\) with \(n=3\) tail \(<10^{-15}\).)

All numbers below are obtained from (5.9)--(5.10) by \textbf{interval arithmetic on ten decimals}: each intermediate quantity is enclosed in an interval whose endpoints are ten-decimal rationals, obtained by rounding \textbf{outwards} the exact rational bounds produced by \(+,-,\times,\div\), and, for square roots, by the elementary rule ``\(a\le\sqrt X\le b\) whenever \(a^2\le X\) and \(b^2\ge X\)'' with ten-decimal rationals \(a,b\). Every single step is a finite rational computation, and every displayed interval below can be re-verified by these four operations alone.

From (5.5), (5.6), (5.4), (5.7) and (5.9)--(5.10): \[\ln22830\in(10.0358307358,\ 10.0358307422),\qquad
\ln22865000\in(16.9451179106,\ 16.9451179168), \tag{5.11}\] \[L_{(1)}\in(1.8817182629,\ 1.8817182642),\qquad
L_{(2)}\in(2.2593490547,\ 2.2593490556), \tag{5.12}\] \[
\begin{aligned}
\ln u&\in(0.4213108974,\ 0.4213108982),\\
\ln(4u)&\in(1.8076052584,\ 1.8076052592),\\
\ln(5u)&\in(2.0307488094,\ 2.0307488108).
\end{aligned}\tag{5.13}
\] (For example, the lower end of (5.11) is \(9\cdot0.6931471804+7\cdot1.0986122884-2\cdot1.9459101492-2\cdot0.0004798041
=6.2383246236+7.6902860188-3.8918202984-0.0009596082=10.0358307358\), and the upper end is \(9\cdot0.6931471806+7\cdot1.0986122888-2\cdot1.9459101484-2\cdot0.0004798040
=10.0358307422\); then \(L_{(1)}=\frac3{16}\ln22830\) gives (5.12) after outward rounding --- e.g.~\(\frac3{16}\cdot10.0358307358=1.88171826296\underline{25}\), which rounds down to \(1.8817182629\). Likewise \(\ln u=\ln3-\ln2+2\operatorname{artanh}\frac{3993}{503993}\) has lower end \(1.0986122884-0.6931471806+2\cdot0.0079228948=0.4213108974\).)

Finally we record the bound for the shift term used twice below: if \(|\delta|\le x\) with \(0<x\le\frac1{10}\), then by (5.8) with \(n=1\) and oddness of \(\operatorname{artanh}\), \[\big|\operatorname{artanh}\delta\big|\le\operatorname{artanh}x
\le x+\frac{x^3}{3(1-x^2)} . \tag{5.14}\]

\subsection{Range I: \(2\le p\le\frac{16}{3}\)}\label{subsec:range1}

Here \(K=0\) and \(\hat c=\frac3{10}\ln22830\), so by (5.11) \[\hat c\in(3.0107492207,\ 3.0107492227) .\] In particular \(\hat c>2\). Since \(X\mapsto X(X-2)\) is increasing for \(X>1\), \[g^2=\hat c(\hat c-2)\in\big(3.0107492207\cdot1.0107492207,\
3.0107492227\cdot1.0107492227\big)\subset(3.0431124285,\ 3.0431124366),\] and, since \(1.7444518991^2\le3.0431124285\) while \(1.7444519015^2\ge3.0431124366\), \[g\in(1.7444518991,\ 1.7444519015). \tag{5.15}\] Hence \(t^*=g/\hat c\in\big(\frac{1.7444518991}{3.0107492227},\frac{1.7444519015}{3.0107492207}\big)
\subset(0.5794079048,\ 0.5794079061)\), which lies in \((0,\frac58)\); so Lemma 5.2 gives the maximum of \(S\) on the \(t\)-interval \([0,\frac58]\) at \(t^*\). With \(\beta=\frac35\), the map \(t\mapsto\frac{t-3/5}{1-3t/5}\) is increasing, so \[\delta\in(-0.0315657689,\ -0.0315657667),\qquad |\delta|\le x:=0.0315657689 .\] By (5.14): \(x^2\le0.0009963978\), \(x^3\le0.0000314521\), and \(\frac{x^3}{3(1-x^2)}\le\frac{0.0000314521}{3\cdot0.9990036022}\le0.0000104945\), so \(\operatorname{artanh}x\le0.0315657689+0.0000104945\le0.0315762634\) and \[2\operatorname{artanh}\delta\ \ge\ -0.0631525268 .\] Therefore, by (5.13), \[\ln(4u)+2\operatorname{artanh}\delta\ \ge\ 1.8076052584-0.0631525268=1.7444527316 ,\] while (recall \(K=0\)) \[K+g\ =\ g\ \le\ 1.7444519015\ <\ 1.7444527316 .\] This is (5.3), so \(\|B\|_p\le u(p-1)\) for \(2\le p\le\frac{16}3\). The margin is \(8.3\cdot10^{-7}\).

\subsection{Range II: \(\frac{16}{3}\le p\le\frac{15}{2}\)}\label{subsec:range2}

By (5.12), \[L_{(2)}-L_{(1)}\in(2.2593490547-1.8817182642,\ 2.2593490556-1.8817182629)
\subset(0.3776307905,\ 0.3776307927),\] \[\hat c=\tfrac{120}{13}\big(L_{(2)}-L_{(1)}\big)\in(3.4858226815,\ 3.4858227019),\] \[
\begin{aligned}
K&=L_{(1)}-\tfrac58\hat c\\
 &\in\bigl(1.8817182629-0.625\cdot3.4858227019,\\
 &\hspace{4.2em}1.8817182642-0.625\cdot3.4858226815\bigr)\\
 &\subset(-0.2969209258,\ -0.2969209117).
\end{aligned}
\] Again \(\hat c>2\), and \[g^2=\hat c(\hat c-2)\in\big(3.4858226815\cdot1.4858226815,\
3.4858227019\cdot1.4858227019\big)\subset(5.1793144038,\ 5.1793145053),\] so, since \(2.2758107135^2\le5.1793144038\) and \(2.2758107359^2\ge5.1793145053\), \[g\in(2.2758107135,\ 2.2758107359).\] Hence \(t^*=g/\hat c\in\big(\frac{2.2758107135}{3.4858227019},\frac{2.2758107359}{3.4858226815}\big)
\subset(0.6528762097,\ 0.6528762201)\), which lies in \((\frac58,\frac{11}{15})
=(0.625,\,0.7333\ldots)\); so Lemma 5.2 applies on the \(t\)-interval \([\frac58,\frac{11}{15}]\). With \(\beta=\frac23\), the increasing map \(t\mapsto\frac{t-2/3}{1-2t/3}\) gives \[\delta\in(-0.0244187284,\ -0.0244187095),\qquad |\delta|\le x:=0.0244187284 .\] By (5.14): \(x^2\le0.0005962743\), \(x^3\le0.0000145603\), and \(\frac{x^3}{3(1-x^2)}\le\frac{0.0000145603}{3\cdot0.9994037257}\le0.0000048564\), so \(\operatorname{artanh}x\le0.0244187284+0.0000048564=0.0244235848\) and \[2\operatorname{artanh}\delta\ \ge\ -0.0488471696 .\] Therefore, by (5.13), \[\ln(5u)+2\operatorname{artanh}\delta\ \ge\ 2.0307488094-0.0488471696=1.9819016398,\] while \[K+g\ \le\ -0.2969209117+2.2758107359=1.9788898242\ <\ 1.9819016398 .\] This is (5.3) with \(\beta=\frac23\) (whose \(2\operatorname{artanh}\beta=\ln5\), so \(\ln u+2\operatorname{artanh}\beta=\ln(5u)\)), so \(\|B\|_p\le u(p-1)\) for \(\frac{16}3\le p\le\frac{15}2\). The margin is \(3.01\cdot10^{-3}\).

\subsection{Completion of the proof}\label{subsec:completion}

Combining the three ranges proves
\[\|B\|_{L^p\to L^p}\le1.523958\,(p-1),\qquad p\ge2.\]
For \(1<p<2\), transposition gives \(\|B\|_p=\|B\|_{p'}\), while
\(p^*-1=p'-1\).  This completes the proof of the main theorem.

\section{Concluding remarks}\label{sec:remarks}

\begin{enumerate}[label=\textup{(\arabic*)},leftmargin=*]
\item
The binding interpolation interval is \([2,16/3]\).  High-precision evaluation
of the explicit interpolation envelopes gives
\[
 \sup_{2\le p\le16/3}\frac{\|B\|_p}{p-1}
 \le 1.5239567163458625\ldots,
\]
with the maximum at \(p=4.7552011223\ldots\).  The other two envelopes are
\(1.5193749633\ldots\) and \(\sqrt{30/13}=1.5191090506\ldots\).  These decimal
values are included only for orientation; the proof of the theorem uses the
outward-rounded rational interval certificate in Section~5.

\item
The condition \(b_2\ge p^2a_0\) in Remark~2.10 is merely a necessary local
condition for the polynomial drift determinant near \(t=0\).  The two profiles
used here choose equality, but no claim of global optimality within all polynomial
profiles, or within all Bellman functions satisfying the drift criterion, is made.

\item
The parameter \(\theta<1\) in the growth axiom is essential at
\(p=16/3\): the monomial \(r^5s^{1/3}\) produces a second derivative of order
\(s^{-5/3}\), which remains locally integrable in the two-dimensional \(y\)
variable.  Compactly supported convolution then preserves the drift inequality
by a direct distributional argument.

\item
The proof identifies two possible sources of further improvement: a stronger
pointwise use of the joint geometry of the radial and tangential components, or
new endpoint Bellman profiles outside the present generalized-polynomial family.
The current argument establishes only the explicit coefficient stated in the main
theorem and does not assert sharpness of the method.
\end{enumerate}

\begin{thebibliography}{99}

\bibitem{BanuelosJanakiraman08}
R. Bañuelos and P. Janakiraman,
\emph{\(L^p\)-bounds for the Beurling--Ahlfors transform},
Trans. Amer. Math. Soc. \textbf{360} (2008), no.~7, 3603--3612.

\bibitem{BanuelosMendez03}
R. Bañuelos and P. J. Méndez-Hernández,
\emph{Space-time Brownian motion and the Beurling--Ahlfors transform},
Indiana Univ. Math. J. \textbf{52} (2003), no.~4, 981--990.

\bibitem{Burkholder84}
D. L. Burkholder,
\emph{Boundary value problems and sharp inequalities for martingale transforms},
Ann. Probab. \textbf{12} (1984), no.~3, 647--702.

\bibitem{Iwaniec82}
T. Iwaniec,
\emph{Extremal inequalities in Sobolev spaces and quasiconformal mappings},
Z. Anal. Anwendungen \textbf{1} (1982), no.~6, 1--16.

\end{thebibliography}

\end{document}
